%% %%%%%%%%%%%%%%%%%%%%%%%%%%%%%%%%%%%%%%%%%%%%%%%%%%%%%%%%%%%%%%%%%%%%%%%%%%%%%%%%%
%% Template for Research proposals for Computer Vision subject
%% %%%%%%%%%%%%%%%%%%%%%%%%%%%%%%%%%%%%%%%%%%%%%%%%%%%%%%%%%%%%%%%%%%%%%%%%%%%%%%%%%
%% Version 1.0 // February 2025
%% %%%%%%%%%%%%%%%%%%%%%%%%%%%%%%%%%%%%%%%%%%%%%%%%%
%% Eduardo FIDALGO
%%  - eduardo.fidalgo@unileon.es
%% %%%%%%%%%%%%%%%%%%%%%%%%%%%%%%%%%%%%%%%%%%%%%%%%%
\documentclass[11pt]{article}
\usepackage{AV_style}
\usepackage{comment}

\usepackage{lipsum}  %% Package to create dummy text (comment or erase before start)

\usepackage{multirow}
\usepackage{diagbox}
\usepackage{xfrac}
\usepackage{tikz}
\usepackage{caption, subcaption}
\usepackage{lscape}
\usepackage{pgfplots}
\usepackage[export]{adjustbox}
\usepackage{xcolor}

%% ===============================================
%% Setting the line spacing (3 options: only pick one)
% \doublespacing
% \singlespacing
\onehalfspacing
%% ===============================================

\setlength{\droptitle}{-5em} %% Don't touch

% %%%%%%%%%%%%%%%%%%%%%%%%%%%%%%%%%%%%%%%%%%%%%%%%%%%%%%%%%%
% SET THE TITLE
% %%%%%%%%%%%%%%%%%%%%%%%%%%%%%%%%%%%%%%%%%%%%%%%%%%%%%%%%%%
% TITLE:
\title{Research Proposal: Detection and Segmentation of Cut Edges
}

% AUTHORS:
\author{
    Pablo Ruíz Morán \\
    \href{mailto:pruizm01@estudiantes.unileon.es}{\texttt{pruizm01@estudiantes.unileon.es}}
    \and
    David Morán Gorgojo \\
    \href{mailto:dmorag03@estudiantes.unileon.es}{\texttt{dmorag03@estudiantes.unileon.es}}
    \and
    Miguel Sánchez Rodriguez \\
    \href{mailto:msancr11@estudiantes.unileon.es}{\texttt{msancr11@estudiantes.unileon.es}}
    \and
    Jaime Alvarado Fernandez \\
    \href{mailto:jalvaf08@estudiantes.unileon.es}{\texttt{jalvaf08@estudiantes.unileon.es}}
}
    
% DATE:
\date{\today}

% %%%%%%%%%%%%%%%%%%%%%%%%%%%%%%%%%%%%%%%%%%%%%%%%%%%%%%%%%%
% %%%%%%%%%%%%%%%%%%%%%%%%%%%%%%%%%%%%%%%%%%%%%%%%%%%%%%%%%%
\begin{document}
% %%%%%%%%%%%%%%%%%%%%%%%%%%%%%%%%%%%%%%%%%%%%%%%%%%%%%%%%%%
% %%%%%%%%%%%%%%%%%%%%%%%%%%%%%%%%%%%%%%%%%%%%%%%%%%%%%%%%%%
% ABSTRACT
% %%%%%%%%%%%%%%%%%%%%%%%%%%%%%%%%%%%%%%%%%%%%%%%%%%%%%%%%%%
% %%%%%%%%%%%%%%%%%%%%%%%%%%%%%%%%%%%%%%%%%%%%%%%%%%%%%%%%%%
{\setstretch{.8}
\maketitle
% %%%%%%%%%%%%%%%%%%
\begin{abstract}
This research project addresses the automated detection and metric estimation of 
cutting edges in rubber strip assembly within tyre manufacturing, developed in the 
context of the \textbf{Michelin Challenge 2} proposed by \textit{Michelin Aranda de 
Duero}. The core challenge consists of accurately localizing cut borders and 
estimating their spatial positions in millimetres from smartphone images captured 
under variable perspectives, illumination, and limited annotated data, where 
precision and robustness are critical for ensuring product quality.

To tackle this problem, we review recent methods (2024--2026) spanning object 
detection, image segmentation, anomaly detection, and classical edge detection, and 
propose four complementary computer vision pipelines. The first is 
\textbf{Grounded-SAM}, which combines Grounding DINO for open-vocabulary detection 
with SAM for precise mask generation, enabling zero-shot segmentation without 
task-specific training. The second is a \textbf{YOLOv9 + EfficientSAM} pipeline, 
which integrates real-time object detection with lightweight promptable segmentation 
for fine-tunable performance on the target dataset. The third is an 
\textbf{anisotropic multi-scale edge detector (MAMDD)}, a training-free classical 
approach that captures thin separation lines under noise and low contrast. The fourth 
is an \textbf{adaptive classical pipeline} based on improved Canny edge detection 
with Otsu thresholding, QR-based homography calibration, and Probabilistic Hough 
Transform, deployable entirely on CPU without any training data.

All four pipelines incorporate QR-code-based spatial calibration to correct 
perspective distortion and provide millimetre-level distance measurements. Together, 
they cover a spectrum from zero-shot deep learning to fully classical approaches, 
offering flexible trade-offs between accuracy, computational cost, and data 
requirements for real industrial deployment.\\

\textit{\textbf{Keywords: } Cut Edge Detection, Object Detection, Image Segmentation, 
Grounded-SAM, YOLOv9, Classical Edge Detection, Industrial Inspection, Michelin 
Challenge}

\end{abstract}
}

% %%%%%%%%%%%%%%%%%%%%%%%%%%%%%%%%%%%%%%%%%%%%%%%%%%%%%%%%%%
% %%%%%%%%%%%%%%%%%%%%%%%%%%%%%%%%%%%%%%%%%%%%%%%%%%%%%%%%%%
% BODY OF THE DOCUMENT
% %%%%%%%%%%%%%%%%%%%%%%%%%%%%%%%%%%%%%%%%%%%%%%%%%%%%%%%%%%
% %%%%%%%%%%%%%%%%%%%%%%%%%%%%%%%%%%%%%%%%%%%%%%%%%%%%%%%%%%

% --------------------
\section{Introduction}
\label{sec:introduction}
% Insert the content of the introduction here.
%    What is the problem?
%    Why is the problem significant or important?
%    A brief review of the state of the art.
%    What are the benefits or importance of this research?
%    What will be presented in this document?

Automatic detection and precise localization of cutting edges in industrial images 
is a key problem in modern manufacturing, driven by the increasing demand for 
automated quality control systems. In tyre assembly processes, rubber strips are 
placed over a cylindrical drum, cut, and subsequently joined at their edges. 
Currently, these tasks are performed manually by workers on the production floor, 
introducing variability, human error, and limiting the scalability of the process. 
Automating this inspection step would directly reduce operational costs and improve 
product consistency across production lines.

Precise cut edge localization is particularly challenging in this setting. Images 
are captured using smartphones by different workers from variable positions and 
angles, requiring robustness to perspective distortion, illumination changes, and 
varying camera distance. The rubber strip and its cut edges often present low 
contrast with respect to the background, and an additional requirement is to provide 
metric measurements in millimetres rather than pixel coordinates, which necessitates 
a reliable pixel-to-metric calibration procedure based on QR codes placed on the 
inspection surface.

These challenges sit at the intersection of several active research directions in 
Deep Learning and Computer Vision. Recent advances have produced powerful methods 
for object detection and segmentation: convolutional approaches such as 
YOLOv8~\cite{Jocher2023YOLOv8} offer real-time detection, while transformer-based 
detectors such as RT-DETRv4~\cite{LiaoRTDETRv4} achieve state-of-the-art results 
on standard benchmarks. Vision-language models such as Grounding 
DINO~\cite{RenGroundedSAM2024} extend detection to open-vocabulary settings without 
task-specific retraining, and the Segment Anything Model 
(SAM)~\cite{KirillovSAM2023} enables promptable pixel-level segmentation for 
arbitrary objects. However, most of these methods are evaluated on general-purpose 
benchmarks, and their applicability to industrial inspection with limited labeled 
data, domain shift, and strict metric accuracy requirements remains an open 
challenge.

The specific objectives of this work are: (i) to accurately detect and segment the 
cut edges of rubber strips from unconstrained smartphone images; (ii) to estimate 
their spatial positions and mutual distances in millimetres using QR-based geometric 
calibration; and (iii) to evaluate the suitability of both deep learning and 
classical approaches under the data and deployment constraints of a real industrial 
environment. To this end, we review twenty recent methods (2024--2025) and propose 
four complementary pipelines: Grounded-SAM for zero-shot open-vocabulary 
segmentation, YOLOv9~+~EfficientSAM for fine-tunable detection and segmentation, 
an anisotropic multi-scale edge detector (MAMDD) for training-free classical edge 
localization, and an adaptive Canny-based pipeline with QR-guided homography 
calibration for fully CPU-deployable metric estimation.

The rest of the document is organized as follows. 
Section~\ref{sec:literatureReview} presents the literature review. 
Section~\ref{sec:selectedMethod} describes the four proposed pipelines. 
Section~\ref{sec:datasets} presents the datasets, Section~\ref{sec:metrics} 
defines the evaluation metrics, Section~\ref{sec:python} covers the available 
Python implementations, and Section~\ref{sec:credits} lists individual 
contributions.
% --------------------
% --------------------
% --------------------
\section{Literature Review}
% --------------------
\label{sec:literatureReview}

The following section reviews recent methods (2024--2025) relevant to the
\textbf{Michelin Challenge 2}, organized into four groups: deep learning-based
object detection and segmentation, anomaly detection, classical computer vision,
and camera calibration. The goal of this review is to identify the most suitable
approaches for cut edge detection and metric estimation under the constraints of
the challenge: limited annotated data, variable acquisition conditions, and
the need for millimetre-level accuracy without specialized hardware.

Among the reviewed deep learning methods, the strongest candidates for this
problem are \textbf{Grounded-SAM}~\cite{RenGroundedSAM2024}, which achieves
$48.7$ mAP in zero-shot settings on 25 diverse real-world datasets without
task-specific training; \textbf{YOLOv9 + EfficientSAM}, combining $55.6\%$
AP$_{50:95}$ detection with lightweight promptable segmentation at 47 img/s;
and \textbf{AnomalyDINO}~\cite{Damm2024AnomalyDINO}, which reaches $96.6\%$
AUROC in one-shot anomaly detection without any fine-tuning. On the classical
side, the \textbf{MAMDD} edge detector~\cite{Liang2025} and the
\textbf{improved Canny pipeline}~\cite{Liu2024b} stand out for their
robustness to noise and direct applicability without training data.

% --------------------
\subsection{Deep Learning: Object Detection and Segmentation}
% --------------------

Jocher et al.~\cite{Jocher2023YOLOv8} present YOLOv8, a next-generation object
detection model developed by Ultralytics. The architecture improves previous YOLO
versions by integrating a CSPNet backbone for enhanced feature extraction, an
FPN+PAN neck for multi-scale feature aggregation, and an anchor-free detection
head that simplifies bounding box prediction and reduces computational complexity.
The model also introduces improved training strategies such as advanced data
augmentation techniques (mosaic and mixup), focal loss for classification, and
mixed-precision training. Benchmark evaluations on Microsoft COCO and Roboflow
100 demonstrate that YOLOv8 achieves $71.5\%$ mAP@0.5 with the YOLOv8x variant
at 11.5 ms inference time.\\

Liao et al.~\cite{LiaoRTDETRv4} propose RT-DETRv4, a real-time object detection
framework that improves DETR-based detectors by distilling semantic knowledge from
Vision Foundation Models into lightweight architectures. The method introduces a
Deep Semantic Injector (DSI) module and a Gradient-guided Adaptive Modulation
(GAM) strategy that dynamically balances semantic transfer during training.
Evaluated on COCO 2017, RT-DETRv4 achieves AP scores of 49.7, 53.5, 55.4, and
57.0 for the S, M, L, and X variants at speeds of 273, 169, 124, and 78 FPS
respectively.\\

Kienzle et al.~\cite{KienzleSegformer2024} propose Segformer++, an efficient
transformer-based architecture for high-resolution semantic segmentation that
adapts token-merging strategies to the Segformer framework without requiring model
retraining. Two configurations are introduced: Segformer++HQ, oriented towards
quality preservation, and Segformer++fast, oriented towards speed. Evaluated on
Cityscapes and ADE20K, Segformer++HQ achieves a mIoU of $82.31$ on Cityscapes
with a $1.61\times$ inference speedup, while Segformer++fast reaches a
$1.94\times$ speedup with only a marginal performance drop.\\

Cavagnero et al.~\cite{CavagneroPEM2024} propose PEM (Prototype-based Efficient
MaskFormer), a transformer-based architecture for image segmentation that
introduces a prototype-based cross-attention mechanism to reduce computational
burden, together with an efficient multi-scale feature pyramid network combining
deformable convolutions and context-based self-modulation. Evaluated on
Cityscapes and ADE20K, PEM achieves $79.9$ mIoU at 13.8 FPS for semantic
segmentation and $61.1$ PQ at 13.8 FPS for panoptic segmentation.\\

Hou et al.~\cite{HouRelationDETR2024} propose Relation-DETR, a transformer-based
object detector that introduces an explicit position relation prior to improve
convergence speed and detection accuracy in DETR architectures. The method
incorporates a position relation encoder that models pairwise interactions between
bounding boxes and uses this information as an attention bias for progressive
attention refinement. Evaluated on COCO 2017, Relation-DETR achieves $51.7$ AP
with ResNet-50, outperforming DINO by $+2.0$ AP and reaching over 40 AP after
only 2 training epochs.\\

Heidari et al.~\cite{HeidariViLUNet2025} introduce ViLU-Net, a hybrid
encoder--decoder architecture that integrates Vision-LSTM (xLSTM) blocks within
a U-Net structure to model long-range spatial dependencies. The model combines CNN
layers for local feature extraction with xLSTM modules for global contextual
information. Evaluated on a CT dataset of retroperitoneal tumours and the FLARE
dataset, ViLU-Net achieves a DSC of $0.9309$ and IoU of $0.8720$, outperforming
methods such as nnU-Net, SwinUNETR, and U-Mamba.\\

Ren et al.~\cite{RenGroundedSAM2024} introduce Grounded-SAM, a modular framework
that combines Grounding DINO, an open-set detector capable of identifying objects
from free-form text prompts, with the Segment Anything Model (SAM), which
generates precise pixel-level masks. Grounding DINO detects candidate regions and
outputs bounding boxes, which are then used as prompts by SAM to produce
segmentation masks. Evaluated on the Segmentation in the Wild (SegInW) benchmark
--- comprising 25 diverse real-world datasets --- the method achieves $48.7$ mAP
in zero-shot settings, demonstrating strong generalization for open-world vision
tasks without additional training.\\

Cheng et al.~\cite{ChengYOLOWorld2024} propose YOLO-World, a real-time
open-vocabulary object detector that extends the YOLO architecture by
incorporating linguistic information through a CLIP-based text encoder and a
Re-parameterizable Vision-Language Path Aggregation Network (RepVL-PAN). The
model achieves $35.4$ AP on the LVIS dataset in zero-shot settings at 52 FPS,
demonstrating strong performance in open-world scenarios with objects not
explicitly included in the training set.\\

Wyatt et al.~\cite{WyattAnoDDPM2022} propose an unsupervised anomaly detection
approach for defect localization in ultrasonic non-destructive testing using a
generative diffusion model. Their method relies on AnoDDPM trained exclusively
on defect-free data to reconstruct normal ultrasonic wave propagation patterns.
During inference, defects are detected by computing the difference between the
input image and the reconstructed output, producing anomaly maps that highlight
defective regions. Applied to a Laser Ultrasonic Visualization Testing (LUVT)
dataset for aluminium inspection, the method achieves an AUROC of 0.909,
outperforming VAE ($0.800$) and f-AnoGAN ($0.755$), and reaching Precision
$0.817$ and Recall $0.820$.\\

Ren et al.~\cite{RenGroundingDINO15} introduce Grounding DINO 1.5, an advanced
open-vocabulary detection model that extends the original Grounding DINO framework
by combining visual features with textual prompts through a transformer-based
architecture. Trained on large-scale grounding datasets, the Pro model achieves
$54.3$ AP on COCO and $55.7$ AP on LVIS-minival in zero-shot settings, while an
optimised Edge variant reaches 75.2 FPS on A100 with TensorRT.\\

A summary of the reviewed methods is provided in
Table~\ref{tab1:literatureReview}.

\begin{table}[phtb!]
\begin{center}
\resizebox{\linewidth}{!}{
\begin{tabular}{l c p{4.5cm} c c c c}
\hline
\multirow{2}{*}{Method} & \multirow{2}{*}{Year} & \multirow{2}{*}{Description} &
\multirow{2}{*}{Dataset} & \multirow{2}{*}{Main Metric} &
\multirow{2}{*}{Code} & \multirow{2}{*}{Speed} \\
 & & & & & & \\
\hline
YOLOv8~\cite{Jocher2023YOLOv8} & 2024 &
    Anchor-free CNN detector with CSPNet backbone and FPN+PAN neck &
    COCO, Roboflow 100 & 71.5\% mAP@0.5 & Yes & 11.5 ms \\
RT-DETRv4~\cite{LiaoRTDETRv4} & 2025 &
    DETR detector with VFM distillation via DSI and GAM &
    COCO 2017 & 49.7--57.0 AP & Yes & 78--273 FPS \\
Segformer++~\cite{KienzleSegformer2024} & 2024 &
    Efficient transformer segmentation via token merging &
    Cityscapes, ADE20K & 82.31 mIoU & Yes & $1.61\times$ speedup \\
PEM~\cite{CavagneroPEM2024} & 2024 &
    Prototype-based efficient MaskFormer &
    Cityscapes, ADE20K & 61.1 PQ / 79.9 mIoU & Yes & 13.8--35.7 FPS \\
Relation-DETR~\cite{HouRelationDETR2024} & 2024 &
    DETR with explicit position relation prior &
    COCO 2017 & 51.7 AP & Yes & N/R \\
ViLU-Net~\cite{HeidariViLUNet2025} & 2025 &
    Hybrid U-Net with CNN and xLSTM blocks &
    CT Tumour, FLARE & DSC 0.9309 / IoU 0.8720 & Yes & N/R \\
Grounded-SAM~\cite{RenGroundedSAM2024} & 2024 &
    Open-vocabulary detection + segmentation pipeline &
    SegInW (zero-shot) & 48.7 mAP & Yes & N/R \\
YOLO-World~\cite{ChengYOLOWorld2024} & 2024 &
    Real-time open-vocabulary YOLO with RepVL-PAN &
    LVIS & 35.4 AP & Yes & 52.0 FPS \\
Grounding DINO 1.5~\cite{RenGroundingDINO15} & 2024 &
    Advanced open-vocabulary detector with ViT-L backbone &
    COCO, LVIS & 54.3 AP / 55.7 AP & Yes & 5.5--75.2 FPS \\
AnoDDPM~\cite{WyattAnoDDPM2022} & 2024 &
    Unsupervised anomaly detection with diffusion models &
    LUVT & AUROC 0.909 & Yes & N/R \\
\hline
\end{tabular}
}
\caption{\label{tab1:literatureReview}Summary of the reviewed deep learning
methods for object detection and segmentation. N/R = Not Reported.}
\end{center}
\end{table}



Ravi et al.~\cite{Ravi2024SAM2} present SAM 2 (Segment Anything Model 2), a
foundation model developed by Meta FAIR for promptable visual segmentation in
images and videos. The architecture consists of a hierarchical Hiera image
encoder pretrained with MAE, a prompt encoder accepting points, bounding boxes
or masks, a two-way transformer mask decoder, and a streaming memory module for
segmentation propagation across frames. Trained on the SA-V dataset comprising
50.9K videos and 35.5M masks, SAM 2 achieves a J\&F of $76.6\%$ on MOSE val,
while being $6\times$ faster than its predecessor on image segmentation
benchmarks, obtaining $61.9\%$ mIoU across 23 zero-shot datasets.\\

Xiong et al.~\cite{Xiong2023EfficientSAM} propose EfficientSAM, a lightweight
version of SAM that addresses the high computational cost of the original model.
Their approach, SAMI (SAM-leveraged Masked Image Pretraining), trains lightweight
ViT encoders to reconstruct feature embeddings from SAM's ViT-H encoder,
transferring representational knowledge into compact models. EfficientSAM-S
reduces SAM's parameter count by $\sim20\times$ (25M vs 636M) while achieving
$44.4$ AP on zero-shot instance segmentation on COCO at 47 images per second,
compared to SAM's 2 images per second.\\

Wang et al.~\cite{Wang2024YOLOv9} propose YOLOv9, a real-time object detection
model that addresses the information bottleneck problem through two contributions:
Programmable Gradient Information (PGI), which preserves complete input
information during training without adding inference cost, and the Generalized
Efficient Layer Aggregation Network (GELAN), a lightweight architecture based on
CSPNet and ELAN. Evaluated on MS COCO, YOLOv9-E achieves $55.6\%$ AP$_{50:95}$
with 57.3M parameters, reducing parameter count by $49\%$ and computation by
$43\%$ compared to YOLOv8-X while improving AP by $0.6\%$.\\

Shi et al.~\cite{Shi2025CarbonBlock} propose YOLOv8-HDSA, an improved
YOLOv8-based instance segmentation algorithm for industrial carbon block
recognition. The architecture introduces a Selective Reinforcement Feature Fusion
Module (SRFF) using Hadamard product and dilated convolution to suppress
background noise, a convolutional self-attention mechanism for fine-grained edge
detail extraction, and Focaler-IoU as a loss function to dynamically adjust
sample weights. Evaluated on a real industrial dataset, YOLOv8-HDSA improved
recognition accuracy by $7.2\%$ and segmentation accuracy by $3.8\%$ over the
baseline YOLOv8.\\

Damm et al.~\cite{Damm2024AnomalyDINO} propose AnomalyDINO, a training-free
framework for one-shot and few-shot industrial anomaly detection. The method
uses DINOv2 as a patch-level feature extractor and computes distances between
patch embeddings of test images and normal reference images to produce anomaly
scores and pixel-level segmentation maps, without requiring fine-tuning or
additional training data. Evaluated on MVTec-AD, AnomalyDINO achieves an AUROC
of $96.6\%$ in the one-shot setting, improving over the previous state of the
art by $3.5\%$.\\

Jiang et al.~\cite{Jiang2025CLIPDINO} propose a zero-shot industrial anomaly
detection framework that fuses CLIP's global semantic embeddings with DINOv2's
multi-scale structural features through a Dual-Modality Attention (DMA) mechanism.
A complementary Stabilized Attention-based Pooling (SAP) module suppresses normal
background dilution using self-generated anomaly heatmaps. Evaluated across seven
industrial benchmarks, the framework achieves an average image-level AUROC of
$93.4\%$, AP of $94.3\%$, pixel-level AUROC of $96.9\%$, and AUPRO of $92.4\%$,
outperforming WinCLIP, AnomalyCLIP, and Crane.\\

Xiong et al.~\cite{Xiong2026SAM2UNet} propose SAM2-UNet, a U-shaped segmentation
framework that leverages the hierarchical Hiera backbone of SAM2 as encoder.
Lightweight adapters with a bottleneck MLP design are inserted before each Hiera
block for parameter-efficient fine-tuning while keeping the encoder frozen. The
decoder follows the classic U-Net design with deep supervision trained with
weighted IoU and binary cross-entropy losses. Evaluated across 18 datasets on five
benchmarks, SAM2-UNet achieves a mDice of $92.8\%$ on Kvasir-SEG and an
S-measure improvement of $4.8\%$ over FEDER on CAMO.\\

Shi et al.~\cite{Shi2024SemiSupervised} propose a semi-supervised segmentation
framework for industrial surface defect detection based on a Two-Branch
Perturbation Cross Pseudo Supervision model (TPcps). One branch applies
image-level perturbations while the second perturbs the feature space via random
channel dropout, with both branches generating mutual pseudo-labels. The network
combines a lightweight MobileNetv2 backbone with DeepLabv3Plus and a dynamic
threshold strategy. Evaluated on KolektorSDD and five self-built datasets, the
method achieves a mIoU $4.39\%$ higher than the state of the art at 64.59 FPS.\\

Liao et al.~\cite{Liao2025CameraCalibration} present a comprehensive survey of
learning-based camera calibration techniques, categorizing existing methods into
four models: pinhole, distortion, cross-view, and cross-sensor. The survey
reviews deep learning strategies for intrinsic and extrinsic parameter estimation,
lens distortion correction, and homography estimation, and compiles a holistic
calibration dataset covering synthetic and real-world images. This work is
directly relevant to the Michelin Challenge, where QR-based calibration from
smartphone images requires robust geometric estimation under perspective
distortion.\\

Heckler-Kram et al.~\cite{HecklerKram2025MVTecAD2} introduce MVTec AD 2, a new
benchmark for unsupervised anomaly detection addressing the saturation of existing
datasets. The benchmark comprises eight scenarios with over 8000 high-resolution
images covering challenging cases including transparent objects, dark-field
illumination, high variance in normal data, and extremely small defects.
Comprehensive evaluation shows that state-of-the-art methods remain below $60\%$
average AU-PRO, demonstrating the benchmark's discriminatory power.\\

A summary of the reviewed methods is provided in
Table~\ref{tab2:literatureReview}.

\begin{table}[phtb!]
\begin{center}
\resizebox{\linewidth}{!}{
\begin{tabular}{l c p{4.5cm} c c c c}
\hline
\multirow{2}{*}{Method} & \multirow{2}{*}{Year} & \multirow{2}{*}{Description} &
\multirow{2}{*}{Dataset} & \multirow{2}{*}{Main Metric} &
\multirow{2}{*}{Code} & \multirow{2}{*}{Speed} \\
 & & & & & & \\
\hline
SAM 2~\cite{Ravi2024SAM2} & 2024 &
    Transformer with streaming memory for promptable segmentation &
    MOSE val, SA-23 & J\&F 76.6\% / mIoU 61.9\% & Yes & 43.8 FPS \\
EfficientSAM~\cite{Xiong2023EfficientSAM} & 2023 &
    Lightweight SAM via SAMI pretraining with ViT-Tiny/Small encoders &
    COCO, LVIS & 44.4 AP (zero-shot) & Yes & 47--54 img/s \\
YOLOv9~\cite{Wang2024YOLOv9} & 2024 &
    Real-time detector with PGI and GELAN &
    MS COCO & 55.6\% AP$_{50:95}$ & Yes & N/R \\
YOLOv8-HDSA~\cite{Shi2025CarbonBlock} & 2025 &
    Improved YOLOv8 with SRFF and Focaler-IoU for industrial segmentation &
    Industrial carbon block & +7.2\% accuracy over YOLOv8 & Yes & N/R \\
AnomalyDINO~\cite{Damm2024AnomalyDINO} & 2024 &
    Training-free few-shot anomaly detection via DINOv2 nearest neighbor &
    MVTec-AD & AUROC 96.6\% (1-shot) & Yes & N/R \\
CLIP-DINOv2~\cite{Jiang2025CLIPDINO} & 2025 &
    Zero-shot anomaly detection with DMA fusion of CLIP and DINOv2 &
    MVTec AD, VisA & AUROC 93.4\% / AP 94.3\% & No & 23 FPS \\
SAM2-UNet~\cite{Xiong2026SAM2UNet} & 2026 &
    U-shaped segmentation using SAM2 Hiera encoder with PEFT adapters &
    18 datasets (5 benchmarks) & mDice 92.8\% (Kvasir) & Yes & N/R \\
Semi-sup. defect seg.~\cite{Shi2024SemiSupervised} & 2024 &
    Two-branch cross pseudo-supervision with DeepLabv3Plus &
    KolektorSDD, Steel & mIoU +4.39\% over SOTA & No & 64.59 FPS \\
Camera Calib. Survey~\cite{Liao2025CameraCalibration} & 2025 &
    Survey of DL-based camera calibration methods &
    Holistic calib. dataset & N/R (survey) & Yes & N/R \\
MVTec AD 2~\cite{HecklerKram2025MVTecAD2} & 2025 &
    Industrial anomaly detection benchmark with 8 challenging scenarios &
    MVTec AD 2 & AU-PRO $<$60\% (SOTA) & No & N/R \\
\hline
\end{tabular}
}
\caption{\label{tab2:literatureReview}Summary of the reviewed methods for
anomaly detection and segmentation foundations. N/R = Not Reported.}
\end{center}
\end{table}

% --------------------
\subsection{Classical Computer Vision}
% --------------------

Yesmin et al.~\cite{Yesmin2025ModifiedWatershed} propose a classical image
segmentation method for liver cancer detection based on a modified
marker-controlled watershed algorithm combined with Sobel edge detection. The
approach applies preprocessing steps including grayscale conversion, gradient
magnitude computation, and morphological operations to extract foreground and
background markers, reducing over-segmentation and noise. Evaluated on MRI images
from the LiTS17 dataset, the method shows improvements in PSNR, SNR, and MSE,
although its reliance on handcrafted preprocessing limits robustness across
diverse imaging conditions.\\

Liang et al.~\cite{Liang2025} propose a noise-robust edge detection method based
on multi-scale anisotropic morphological Gaussian kernels, designed to overcome
limitations of classical detectors such as Canny in noisy environments. The
approach introduces an AMDD framework to capture multi-directional intensity
variations and combines edge strength and direction maps to improve localization
and robustness. Evaluated on BSDS500, NYUDv2, and PASCAL VOC, the method achieves
competitive precision--recall scores, outperforming classical detectors while
remaining comparable to state-of-the-art techniques.\\

Sen et al.~\cite{Sen2025FeatureEngineering} propose a classical image
classification framework based on the fusion of handcrafted features, combining
permutation entropy (PE), Histogram of Oriented Gradients (HOG), and Local Binary
Patterns (LBP) with an SVM classifier. The method achieves competitive performance
on Fashion-MNIST, KMNIST, EMNIST, and CIFAR-10, offering improved interpretability
and computational efficiency compared to deep learning approaches.\\

Xu~\cite{Xu2025IndustrialInspection} presents an overview of automated image
processing for industrial inspection, combining classical techniques such as noise
filtering, histogram equalization, adaptive thresholding, and edge detection with
machine learning-based classification using GLCM and LBP descriptors. While the
approach is adaptable to complex industrial conditions, it relies heavily on system
design and parameter tuning, which may limit generalization.\\

Song and Wen~\cite{Song2025BezierHarris} propose a 3D modeling method that
combines Bezier curves with an improved Harris corner detector. The pipeline
includes denoising, feature extraction using an enhanced CSS-Harris algorithm,
contour generation via Bezier curves, and surface reconstruction with Delaunay
triangulation, showing strong performance in accuracy and speed for geometric
modeling tasks.\\

Hussain~\cite{Hussain2025FuzzySegmentation} proposes a region-based image
segmentation method that combines fuzzy logic with geometric features using an
energy-based level-set framework to balance region homogeneity, edge preservation,
and boundary smoothness. The method achieves $96.8\%$ accuracy and F1 of $95.1\%$,
although it requires parameter tuning and may face scalability challenges in
real-time applications.\\

Almohimeed et al.~\cite{Almohimeed2024SPORGRBVS} propose a retinal vessel
segmentation method combining region growing with the sandpiper optimization
algorithm (SPO) for automatic seed point selection and thresholding. The approach
achieves up to $98.7\%$ accuracy on DRIVE, STARE, and CHASE\_DB1, although it
increases computational complexity and requires parameter tuning.\\

Jia et al.~\cite{Jia2025RidgeSegmentation} propose an unsupervised
ridge-patterned instance segmentation method for high-resolution 3D point clouds,
following a four-stage pipeline of ridge extraction, enhancement, adaptive region
growing, and plane fitting. The method achieves up to $95.5$ mIoU and is
resolution-independent without requiring training data or GPU, although it may
struggle with smooth surfaces and complex thin structures.\\

Guo et al.~\cite{Guo2024SubpixelFlange} propose a machine vision method for
sub-pixel accurate measurement of flange disk dimensions, combining an improved
Canny edge detector with a refined Zernike moment algorithm for sub-pixel
refinement followed by least squares fitting. The method achieves reduced
measurement errors compared to traditional approaches, although it increases
computational complexity and requires careful parameter tuning.\\

Zangana et al.~\cite{Zangana2024EdgeReview} present a comprehensive review of
edge detection techniques covering classical methods (Sobel, Canny, Prewitt) and
advanced approaches based on deep learning, fuzzy logic, and optimization. The
study identifies key challenges including noise sensitivity, computational cost,
and limited generalization, pointing towards future research in real-time and
explainable edge detection systems.\\

A summary of the reviewed methods is provided in
Table~\ref{tab3:literatureReview}.

\begin{table}[phtb!]
\begin{center}
\resizebox{\linewidth}{!}{
\begin{tabular}{l c p{4.5cm} c c c c}
\hline
\multirow{2}{*}{Method} & \multirow{2}{*}{Year} & \multirow{2}{*}{Description} &
\multirow{2}{*}{Dataset} & \multirow{2}{*}{Main Metric} &
\multirow{2}{*}{Code} & \multirow{2}{*}{Speed} \\
 & & & & & & \\
\hline
Modified Watershed~\cite{Yesmin2025ModifiedWatershed} & 2025 &
    Marker-controlled watershed for segmentation and noise removal &
    LiTS17 & PSNR, SNR, MSE (improved) & Yes & N/R \\
MAMDD~\cite{Liang2025} & 2025 &
    Multi-scale anisotropic morphological Gaussian kernel edge detection &
    BSDS500, NYUDv2, PASCAL VOC & PR curves, FOM & No & N/R \\
HOG + LBP + Entropy + SVM~\cite{Sen2025FeatureEngineering} & 2025 &
    Feature fusion for classical image classification &
    CIFAR-10, Fashion-MNIST & 94.81\% accuracy & Yes & N/R \\
Industrial Pipeline~\cite{Xu2025IndustrialInspection} & 2025 &
    Automated vision pipeline for industrial inspection &
    Industrial datasets & N/R & No & N/R \\
Harris + Bezier + Delaunay~\cite{Song2025BezierHarris} & 2025 &
    Corner detection and geometric modeling for 3D reconstruction &
    3D object modeling & 120 pts / 3.2 s & Yes & 3.2 s \\
Fuzzy Segmentation~\cite{Hussain2025FuzzySegmentation} & 2025 &
    Region segmentation using fuzzy logic and geometry &
    Segmentation datasets & 96.8\% accuracy, F1: 95.1\% & Yes & 18.5 s \\
SPORG-RBVS~\cite{Almohimeed2024SPORGRBVS} & 2024 &
    Region growing with optimization for retinal vessel segmentation &
    DRIVE, STARE, CHASE\_DB1 & 98.68\% accuracy & No & N/R \\
Ridge Segmentation~\cite{Jia2025RidgeSegmentation} & 2025 &
    Ridge-pattern segmentation for 3D point clouds &
    ABC-10, OmniObject3D-8 & 95.52 mIoU & No & N/R \\
Subpixel Edge Detection~\cite{Guo2024SubpixelFlange} & 2024 &
    Canny + Zernike subpixel edge detection for industrial measurement &
    Industrial measurement & Reduced absolute/relative error & No & N/R \\
Edge Detection Review~\cite{Zangana2024EdgeReview} & 2024 &
    Review of classical and modern edge detection techniques &
    Literature & Precision, Recall, F1 & No & N/R \\
\hline
\end{tabular}
}
\caption{\label{tab3:literatureReview}Summary of the reviewed classical computer
vision methods. N/R = Not Reported.}
\end{center}
\end{table}


Mira et al.~\cite{Mira2024} present a comparative study between classical and
deep learning approaches for olive fly detection, evaluating combinations such as
HOG+SVM, LeNet-5 CNN, and PCA+Random Forest. Their experiments show that PCA
combined with Random Forest achieves the best performance, reaching up to $99\%$
accuracy and F1-scores around 0.96, while being more efficient for deployment on
edge devices compared to CNN-based solutions.\\

González et al.~\cite{Gonzalez2024} propose a robot-assisted computer vision
system for adaptive edge finishing in industrial components. The method integrates
camera calibration, adaptive binarization, and contour matching using Hu moments
to compute deviations with respect to CAD models, reducing the contour RMSE from
1.050 mm to 0.383 mm and demonstrating high precision in industrial environments.\\

Muzakki et al.~\cite{Muzakki2025} introduce a classical approach for bone
fracture detection in X-ray images using K-Nearest Neighbors combined with HOG,
LBP, and SIFT features. Results on 10,580 images indicate that HOG+KNN achieves
the best performance at $99.03\%$ accuracy.\\

Liu et al.~\cite{Liu2024a} propose a lightweight defect detection framework based
on HOG feature extraction, PCA dimensionality reduction, and a LightGBM classifier
within a grid-based detection approach. The method achieves up to $98\%$ accuracy
with low memory consumption and 0.05 s inference time, making it suitable for
resource-constrained environments.\\

Eryadi~\cite{Eryadi2025} conducts a comparative analysis of crack detection
methods using HOG features combined with SVM and XGBoost classifiers. Experiments
on 40,000 images show that XGBoost achieves the highest accuracy ($99.05\%$) and
faster inference times, while SVM produces fewer false positives.\\

Hung et al.~\cite{Hung2024} propose a tool wear classification system combining
GLCM texture features with SVM and CNN models. CNN models achieve over $93\%$
accuracy, while SVM performance drops to around $70\%$ in complex scenarios,
highlighting the limitations of classical methods for highly variable patterns.\\

Priananda et al.~\cite{Priananda2024} develop a classical pipeline for
phase-resolved partial discharge (PRPD) pattern classification using HOG features,
PCA, and ensemble classifiers. Bagging Trees achieves $89.79\%$ accuracy,
demonstrating the effectiveness of classical pipelines in industrial diagnostic
tasks.\\

Liu et al.~\cite{Liu2024b} propose an improved Canny-based edge detection method
incorporating adaptive filtering, Scharr gradient operators, and Otsu thresholding.
The approach achieves an average precision of up to 0.92 with improved robustness
to noise compared to traditional edge detectors.\\

Ding et al.~\cite{Ding2025Homography} present a closed-form solution for
homography decomposition using geometric constraints derived from three 3D
correspondences. The method achieves accuracy comparable to traditional SVD-based
approaches while being 4--7 times faster and more robust in challenging scenarios.\\

Zhang et al.~\cite{Zhang2024Calibration} propose a camera calibration method
based on a single checkerboard image using a decoupled parameter estimation
strategy, improving calibration accuracy and stability while reducing the
complexity of traditional multi-image calibration techniques.\\

Ge et al.~\cite{Ge2025} propose a defect classification framework combining
improved Canny edge detection, Hu moment features, and an SVM classifier. The
method achieves $97.2\%$ accuracy while maintaining low computational cost,
offering performance comparable to deep learning approaches.\\

A summary of the reviewed methods is provided in
Table~\ref{tab4:literatureReview}.

\begin{table}[phtb!]
\begin{center}
\resizebox{\linewidth}{!}{
\begin{tabular}{l c p{4.5cm} c c c c}
\hline
\multirow{2}{*}{Method} & \multirow{2}{*}{Year} & \multirow{2}{*}{Description} &
\multirow{2}{*}{Dataset} & \multirow{2}{*}{Main Metric} &
\multirow{2}{*}{Code} & \multirow{2}{*}{Speed} \\
 & & & & & & \\
\hline
HOG + SVM / PCA + RF~\cite{Mira2024} & 2024 &
    Classical CV comparison for object detection on edge devices &
    Olive fly dataset & Accuracy: 99\%, F1: 0.96 & No & 9 ms/ROI \\
Adaptive Edge Finishing~\cite{Gonzalez2024} & 2024 &
    Robot vision with adaptive binarization and Hu moment contour matching &
    Aeronautical parts & RMSE: 0.383 mm & No & 31.7 s/image \\
HOG / LBP / SIFT + KNN~\cite{Muzakki2025} & 2025 &
    Classical feature extraction for fracture detection &
    10,580 X-ray images & Accuracy: 99.03\% & No & N/R \\
Grid-HOG-PCA-LightGBM~\cite{Liu2024a} & 2024 &
    Lightweight grid-based defect detection &
    Industrial surface (399 images) & Accuracy: 98\% & Yes & 0.05 s \\
HOG + SVM / XGBoost~\cite{Eryadi2025} & 2025 &
    Crack detection with HOG and ensemble classifiers &
    40,000 crack images & Accuracy: 99.05\% & No & N/R \\
GLCM + SVM / CNN~\cite{Hung2024} & 2024 &
    Tool wear classification with texture features &
    Tool wear dataset & Accuracy: 93\%+ (CNN) & No & N/R \\
HOG + PCA + Ensemble~\cite{Priananda2024} & 2024 &
    PRPD pattern classification with HOG and ensemble classifiers &
    PRPD dataset & Accuracy: 89.79\% & No & N/R \\
Improved Canny~\cite{Liu2024b} & 2024 &
    Adaptive Canny with Scharr gradients and Otsu thresholding &
    Rebar dataset & AP: 0.92 & Yes & N/R \\
Homography Decomp.~\cite{Ding2025Homography} & 2025 &
    Closed-form homography from geometric constraints &
    Benchmark dataset & 4--7$\times$ faster than SVD & Yes & 4--7$\times$ \\
Camera Calibration~\cite{Zhang2024Calibration} & 2024 &
    Single-image calibration with decoupled estimation &
    Checkerboard dataset & Improved accuracy & Yes & N/R \\
Canny + Hu + SVM~\cite{Ge2025} & 2025 &
    Defect classification with edge detection and shape descriptors &
    Automotive flange & Accuracy: 97.2\% & No & N/R \\
\hline
\end{tabular}
}
\caption{\label{tab4:literatureReview}Summary of the reviewed classical computer
vision methods with industrial and calibration focus. N/R = Not Reported.}
\end{center}
\end{table}

% --------------------
\subsection*{Comparative Analysis and Method Selection}
% --------------------

The reviewed methods span a wide spectrum of approaches, from zero-shot
vision-language models to fully classical pipelines. To select the most suitable
methods for the Michelin Challenge 2, four primary criteria were applied:
(i)~\textbf{data requirements} --- whether the method requires large annotated
datasets or can operate in zero-shot or few-shot settings;
(ii)~\textbf{computational cost} --- whether the method requires GPU
infrastructure or can run on standard CPU hardware;
(iii)~\textbf{applicability to industrial edge detection} --- whether the method
has been validated in industrial or low-contrast scenarios; and
(iv)~\textbf{integration with geometric calibration} --- whether the method can
be combined with QR-based metric estimation.

Table~\ref{tab:comparison} summarizes the comparison of the main candidate
methods against these criteria.

\begin{table}[ht]
\centering
\resizebox{\linewidth}{!}{
\begin{tabular}{l c c c c}
\hline
\textbf{Method} & \textbf{Training data} & \textbf{Hardware} &
\textbf{Industrial validation} & \textbf{Calibration-compatible} \\
\hline
Grounded-SAM & Zero-shot & GPU & Partial & Yes \\
YOLOv9 + EfficientSAM & Fine-tuning & GPU & Yes & Yes \\
AnomalyDINO & Few-shot & GPU & Yes (MVTec) & No \\
CLIP-DINOv2 & Zero-shot & GPU & Yes (7 benchmarks) & No \\
MAMDD & None & CPU & Partial & Yes \\
Adaptive Canny pipeline & None & CPU & Yes & Yes \\
\hline
\end{tabular}
}
\caption{Comparison of candidate methods against the main selection criteria
for the Michelin Challenge 2.}
\label{tab:comparison}
\end{table}

Anomaly detection methods such as AnomalyDINO and CLIP-DINOv2, while achieving
strong results on industrial benchmarks, are designed to detect deviations from
normal patterns rather than to localize and measure geometric structures. This
makes them less suitable as primary methods for precise cut edge localization and
metric estimation, although they could complement the pipeline for quality control
purposes.\\

Deep learning methods such as Grounded-SAM and YOLOv9~+~EfficientSAM offer strong
detection and segmentation capabilities but require GPU hardware and, in the case
of YOLOv9, labeled training data for fine-tuning. Their main advantage over
classical approaches is robustness to appearance variability and the ability to
generalize across different acquisition conditions. However, their deployment in a
factory environment with limited infrastructure introduces additional complexity.\\

Classical pipelines based on MAMDD and adaptive Canny edge detection require no
training data or GPU resources and offer full interpretability and ease of
deployment. Their main limitation is sensitivity to parameter tuning and reduced
robustness under large variations in illumination or surface texture compared to
end-to-end trained models.\\

Based on this analysis, four complementary pipelines are proposed in
Section~\ref{sec:selectedMethod}, covering the full spectrum from zero-shot deep
learning to classical CPU-deployable methods, offering flexible trade-offs between
accuracy, data requirements, and deployment constraints.\\
% --------------------
% --------------------
\section{Selected Method}
% --------------------
\label{sec:selectedMethod}
% Each student must describe one method in detail (at least one traditional 
% and one deep learning approach per group).
% TODO: add subsections for the traditional method(s) selected by other group members.

Grounded-SAM \cite{RenGroundedSAM2024} has been selected as the primary deep 
learning due to its ability to combine 
open-vocabulary object detection with high-quality segmentation in a unified pipeline.

% -----------------------------------------------
\subsection{Grounded-SAM (Deep Learning Approach)}
% -----------------------------------------------

\subsubsection{Description and motivation}

Grounded-SAM \cite{RenGroundedSAM2024} is a computer vision framework that combines 
Grounding DINO, an open-vocabulary object detector, with the Segment Anything Model (SAM) 
to perform automatic object detection and segmentation. The method enables detecting objects 
using natural language prompts and generating precise pixel-level segmentation masks, without 
requiring task-specific retraining.

This method was selected for the Michelin Challenge 2 because the combination of detection 
and segmentation in a single pipeline allows the system to:
\begin{itemize}
    \item detect QR codes for spatial calibration,
    \item segment the table or conveyor belt area,
    \item detect and segment rubber strip cut edges,
    \item estimate distances between detected elements using geometric analysis.
\end{itemize}

\subsubsection{Architecture}

Grounded-SAM has a modular design composed of two main deep learning models working sequentially.

\paragraph{Grounding DINO (Object Detection Stage).}
Grounding DINO is a transformer-based object detector that performs open-vocabulary 
detection using natural language prompts. The transformer aligns visual features with text 
embeddings through cross-attention mechanisms, allowing the model to associate image 
regions with semantic concepts. Its main components are:

\begin{itemize}
    \item \textbf{Backbone network:} Extracts multi-scale visual features from the input image.
    \item \textbf{Text encoder:} Encodes the input text prompt (e.g., ``QR code'' or 
    ``rubber strip edge'') into text embeddings.
    \item \textbf{Transformer encoder-decoder:} Processes visual features and aligns them 
    with text embeddings through cross-attention layers.
    \item \textbf{Detection head:} Predicts bounding boxes and confidence scores 
    corresponding to the detected objects.
\end{itemize}

The output of this stage is a set of bounding boxes associated with the objects described 
in the text prompt.

\paragraph{Segment Anything Model (SAM) (Segmentation Stage).}
SAM performs high-quality segmentation based on input prompts. Its architecture includes:

\begin{itemize}
    \item \textbf{Image encoder (Vision Transformer):} Extracts high-level visual feature 
    embeddings from the entire input image. This stage runs only once per image.
    \item \textbf{Prompt encoder:} Encodes prompts such as bounding boxes or points 
    into dense and sparse embeddings.
    \item \textbf{Mask decoder:} Generates the final pixel-level segmentation mask by 
    combining image embeddings with prompt embeddings through cross-attention and 
    upsampling operations.
\end{itemize}

SAM receives the bounding boxes produced by Grounding DINO as spatial prompts and 
produces pixel-level segmentation masks for the detected objects.

\subsubsection{Computer vision pipeline}

\paragraph{Step 1: Image acquisition.}
An RGB image of the inspection area is captured using a smartphone. Since images are 
taken by different workers at variable positions and angles, the system must be robust to 
perspective distortion and changes in illumination.

\paragraph{Step 2: Preprocessing.}
Before feeding the image into the neural networks, the following operations are applied:
\begin{itemize}
    \item \textbf{Image resizing:} The image is resized to match the input resolution 
    expected by each model.
    \item \textbf{Pixel normalization:} Pixel values are normalized to follow the 
    distribution used during training.
    \item \textbf{Tensor conversion:} The image is converted into a PyTorch tensor for 
    neural network processing.
    \item \textbf{QR code detection:} QR codes in the scene are detected and decoded 
    to extract the calibration reference points for metric distance estimation.
\end{itemize}

\paragraph{Step 3: Feature extraction.}
The preprocessed image is processed by the Grounding DINO backbone, which extracts 
multi-scale visual feature maps capturing both global and local visual information from 
the scene.

\paragraph{Step 4: Object detection with Grounding DINO.}
Grounding DINO receives the visual features and a text prompt describing the objects 
of interest. Examples of prompts used in this application include:
\begin{itemize}
    \item \texttt{"QR code"}
    \item \texttt{"rubber strip"}
    \item \texttt{"cutting edge"}
\end{itemize}
The output of this stage is a set of bounding boxes localizing the objects described 
by the prompt.

\paragraph{Step 5: Segmentation with SAM.}
The bounding boxes produced by Grounding DINO are passed as spatial prompts to SAM, 
which generates pixel-level segmentation masks for the detected objects.

\paragraph{Step 6: Edge extraction and post-processing.}
Once segmentation masks are obtained, image processing techniques are applied to 
extract geometric information:
\begin{itemize}
    \item contour detection,
    \item edge extraction,
    \item boundary analysis.
\end{itemize}
These techniques allow identifying the exact borders of the segmented rubber strips 
or cutting edges.

\paragraph{Step 7: Geometric analysis and metric estimation.}
Using the segmented masks and the QR codes as reference markers, the system computes:
\begin{itemize}
    \item distances between cutting edges,
    \item distance from the edge to the table border,
    \item object positions in the coordinate system defined by the QR codes, where 
    point $(0, 0)$ is the centre of the upper-left QR code.
\end{itemize}
The pixel-to-millimetre conversion is derived from the known physical distance between 
QR codes (800 mm legs of a right triangle) using homography or affine transformations 
to correct for perspective distortion.

\subsubsection{Output}

The output of the Grounded-SAM pipeline includes:
\begin{itemize}
    \item bounding boxes of detected objects,
    \item pixel-level segmentation masks of the rubber strip cut edges,
    \item pixel coordinates of detected edge boundaries,
    \item metric measurements of edge positions and distances in millimetres.
\end{itemize}

\subsubsection{Evaluation metrics}

The performance of Grounded-SAM is evaluated using the following metrics. For the object 
detection stage performed by Grounding DINO: Average Precision (AP) and mean Average 
Precision (mAP), which evaluate the accuracy of predicted bounding boxes with respect to 
ground truth annotations. For the segmentation stage performed by SAM: Intersection over 
Union (IoU) and mask IoU, which measure the overlap between the predicted segmentation 
mask and the ground truth mask. These performance metrics are described in detail in
Section~\ref{sec:metrics}

\subsubsection{Applicability to the Michelin Challenge}

Grounded-SAM is particularly suited for the Michelin Challenge 2 because it integrates 
object detection and segmentation in a unified pipeline without requiring task-specific 
training data. In this context, the method can detect QR codes for spatial calibration, 
segment the table or conveyor belt area, identify the cutting edges of rubber strips, and 
compute metric distances between detected elements. The combination of open-vocabulary 
detection and precise segmentation enables accurate localization of relevant structures 
and supports the measurement tasks required in the automated inspection system.

\subsection{YOLOv9 + EfficientSAM (Deep Learning Approach)}

\subsubsection{Description and motivation}

YOLOv9 + EfficientSAM is a two-stage computer vision pipeline that combines a 
real-time object detector with a lightweight segmentation model to perform accurate 
detection and pixel-level segmentation of rubber strip cut edges. YOLOv9 
\cite{WangYOLOv92024} is a state-of-the-art real-time object detector that introduces 
Programmable Gradient Information (PGI) and the Generalized Efficient Layer Aggregation 
Network (GELAN) to address the information bottleneck problem in deep networks. 
EfficientSAM \cite{XiongEfficientSAM2023} is a lightweight variant of the Segment 
Anything Model (SAM) that leverages masked image pretraining to achieve high-quality 
segmentation with significantly reduced computational cost.

This method was selected for the Michelin Challenge 2 because the combination of 
detection and segmentation in a sequential pipeline allows the system to:
\begin{itemize}
    \item detect QR codes for spatial calibration and perspective correction,
    \item segment the table and rubber band area,
    \item detect and segment the cut edges of the rubber strip,
    \item estimate metric distances between detected elements using geometric analysis.
\end{itemize}

\subsubsection{Architecture}

The pipeline is composed of two main deep learning models working sequentially.

\paragraph{YOLOv9 (Object Detection Stage).}
YOLOv9 \cite{WangYOLOv92024} is a real-time object detector built on two core 
contributions:

\begin{itemize}
    \item \textbf{GELAN (Generalized Efficient Layer Aggregation Network):} A lightweight 
    network architecture based on gradient path planning that combines CSPNet and ELAN 
    blocks. GELAN achieves better parameter utilization than depth-wise convolution-based 
    designs while maintaining high inference speed and accuracy.
    \item \textbf{PGI (Programmable Gradient Information):} An auxiliary supervision 
    framework composed of an auxiliary reversible branch and multi-level auxiliary 
    information. PGI ensures that reliable gradient information is available throughout 
    training, preventing information loss caused by the information bottleneck in deep 
    networks. The auxiliary reversible branch is removed at inference time, so it 
    introduces no additional computational cost.
\end{itemize}

The main components of YOLOv9 are:
\begin{itemize}
    \item \textbf{Backbone (GELAN):} Extracts multi-scale visual features from the input 
    image while retaining complete information through gradient path planning.
    \item \textbf{Neck (PAN):} Aggregates features from different scales to improve 
    detection of objects of varying sizes.
    \item \textbf{Detection head:} Predicts bounding boxes and confidence scores for 
    the detected objects.
\end{itemize}

The output of this stage is a set of bounding boxes localizing the cut edges in 
the image.

\paragraph{EfficientSAM (Segmentation Stage).}
EfficientSAM \cite{XiongEfficientSAM2023} is a lightweight variant of SAM obtained 
through a pretraining method called SAMI (SAM-leveraged Masked Image pretraining). 
Instead of reconstructing raw pixel values as in standard masked autoencoders, SAMI 
trains a lightweight Vision Transformer encoder to reconstruct the feature embeddings 
produced by the SAM ViT-H image encoder. This allows the lightweight encoder to acquire 
the strong visual representations of SAM at a fraction of the computational cost. 
Its architecture includes:

\begin{itemize}
    \item \textbf{Lightweight image encoder (ViT-Tiny or ViT-Small):} Pretrained with 
    SAMI to extract high-quality visual feature embeddings from the input image.
    \item \textbf{Prompt encoder:} Encodes spatial prompts such as bounding boxes or 
    points into dense and sparse embeddings.
    \item \textbf{Mask decoder:} Generates the final pixel-level segmentation mask by 
    combining image embeddings with prompt embeddings through cross-attention and 
    upsampling operations.
\end{itemize}

EfficientSAM receives the bounding boxes produced by YOLOv9 as spatial prompts and 
generates pixel-level segmentation masks for the detected cut edges.

\subsubsection{Computer vision pipeline}

\paragraph{Step 1: Image acquisition.}
An RGB image of the inspection area is captured using a smartphone. Since images are 
taken by different workers at variable positions and angles, the system must be robust 
to perspective distortion and illumination changes.

\paragraph{Step 2: Preprocessing.}
Before feeding the image into the neural networks, the following operations are applied:
\begin{itemize}
    \item \textbf{QR code detection and decoding:} The three QR codes present in the 
    scene are detected and decoded using OpenCV. Their pixel coordinates are extracted 
    to define the spatial reference system, where point $(0, 0)$ corresponds to the 
    centre of the upper-left QR code.
    \item \textbf{Homography estimation:} Using the known physical distance between 
    QR codes (800 mm legs of a right triangle) and their detected pixel positions, a 
    homography matrix is computed to correct perspective distortion and derive the 
    pixel-to-millimetre conversion factor.
    \item \textbf{Image resizing and normalization:} The image is resized to match the 
    input resolution expected by each model, pixel values are normalized, and the image 
    is converted into a PyTorch tensor.
\end{itemize}

\paragraph{Step 3: Feature extraction.}
The preprocessed image is processed by the YOLOv9 GELAN backbone, which extracts 
multi-scale visual feature maps capturing both global and local information from 
the scene.

\paragraph{Step 4: Object detection with YOLOv9.}
YOLOv9 processes the extracted features and predicts bounding boxes for the objects 
of interest. In this application, the model is fine-tuned on the Michelin dataset 
to detect:
\begin{itemize}
    \item QR codes (used as calibration markers),
    \item the rubber strip,
    \item the cut edges of the rubber strip.
\end{itemize}
The output of this stage is a set of bounding boxes localizing the cut edges in 
pixel coordinates.

\paragraph{Step 5: Segmentation with EfficientSAM.}
The bounding boxes produced by YOLOv9 are passed as spatial prompts to EfficientSAM, 
which generates pixel-level segmentation masks for the detected cut edges. EfficientSAM 
processes each bounding box independently and outputs a binary mask indicating the 
exact pixel region corresponding to each cut edge.

\paragraph{Step 6: Edge extraction and post-processing.}
Once segmentation masks are obtained, classical image processing techniques are applied 
to extract geometric information:
\begin{itemize}
    \item contour detection,
    \item edge extraction,
    \item boundary analysis.
\end{itemize}
These operations allow identifying the precise boundary coordinates of each segmented 
cut edge.

\paragraph{Step 7: Geometric analysis and metric estimation.}
Using the segmented masks and the homography derived from the QR codes, the system 
computes:
\begin{itemize}
    \item the position of the cut edges in the coordinate system defined by the 
    QR codes, where point $(0, 0)$ is the centre of the upper-left QR code,
    \item the distance from each cut edge to the lower border of the table,
    \item the distance between the two cut edges.
\end{itemize}
All measurements are expressed in millimetres using the pixel-to-millimetre conversion 
factor estimated in the preprocessing step.

\subsubsection{Output}

The output of the YOLOv9 + EfficientSAM pipeline includes:
\begin{itemize}
    \item bounding boxes of the detected cut edges,
    \item pixel-level segmentation masks of the rubber strip cut edges,
    \item pixel coordinates of the detected edge boundaries,
    \item metric measurements of edge positions and distances in millimetres, 
    referenced to the QR code coordinate system.
\end{itemize}

\subsubsection{Evaluation metrics}

The performance of the pipeline is evaluated using the following metrics. For the 
object detection stage performed by YOLOv9: Average Precision (AP) and mean Average 
Precision (mAP), which evaluate the accuracy of predicted bounding boxes with respect 
to ground truth annotations. For the segmentation stage performed by EfficientSAM: 
Intersection over Union (IoU) and mask IoU, which measure the overlap between the 
predicted segmentation mask and the ground truth mask. These performance metrics are 
described in detail in Section~\ref{sec:metrics}.

\subsubsection{Applicability to the Michelin Challenge}

YOLOv9 + EfficientSAM is well suited for the Michelin Challenge 2 because it provides 
a complete and efficient pipeline for detection, segmentation and metric estimation. 
YOLOv9 delivers accurate and fast detection of cut edges even under variable 
illumination and perspective conditions, while EfficientSAM produces precise 
pixel-level segmentation masks at a fraction of the computational cost of the original 
SAM. The use of QR codes as calibration markers and the homography-based perspective 
correction ensure accurate metric estimation regardless of the camera position. 
Together, these components enable reliable localization and measurement of rubber 
strip cut edges in the automated inspection system.


\subsection{Noise-Robust Edge Detection (Classical Approach)}

\subsubsection{Description and motivation}

The selected method is a classical edge detection framework designed to improve robustness against noise and irregular textures by combining multi-scale processing with anisotropic morphological operations. Unlike traditional gradient-based detectors such as Canny or Sobel, this approach introduces a multi-scale anisotropic morphological directional derivative (AMDD) model that captures local intensity variations across multiple orientations.

This method was selected for the Michelin Challenge 2 because the task involves detecting thin separation lines between rubber strips under challenging conditions such as:
\begin{itemize}
    \item low contrast between adjacent regions,
    \item surface texture and noise from rubber material,
    \item varying illumination and acquisition conditions.
\end{itemize}

The proposed method improves edge continuity and robustness, making it suitable for accurately identifying separation boundaries in industrial environments.

\subsubsection{Architecture}

The method follows a multi-stage classical computer vision pipeline composed of the following components:

\textbf{Multi-scale processing.} The input image is processed at multiple scales to capture both fine and coarse edge structures, improving detection robustness across different spatial resolutions.

\textbf{Anisotropic morphological filtering.} Directional morphological Gaussian kernels are applied to enhance edge responses along specific orientations while suppressing noise. This allows better discrimination of relevant edges from background texture.

\textbf{Directional derivative computation (AMDD).} The anisotropic morphological directional derivative is computed to estimate intensity variations across multiple directions, producing:
\begin{itemize}
    \item Edge Strength Map (ESM),
    \item Edge Direction Map (EDM).
\end{itemize}

\textbf{Fusion strategy.} The ESM and EDM are combined to generate a final edge map that improves localization accuracy and continuity.

\textbf{Post-processing.} Additional steps such as contrast equalization and spatial filtering are applied to refine edge quality and suppress residual noise.

\subsubsection{Computer vision pipeline}

\textbf{Step 1: Image acquisition.} An RGB image of the inspection surface is captured using a mobile device or industrial camera. Images may present perspective distortion, varying illumination, and noise.

\textbf{Step 2: Preprocessing.}
\begin{itemize}
    \item Grayscale conversion,
    \item Gaussian smoothing,
    \item Contrast normalization.
\end{itemize}

\textbf{Step 3: Multi-scale feature extraction.} The image is processed at different scales to capture edge structures of varying thickness.

\textbf{Step 4: Edge detection (AMDD).} Directional derivatives are computed using anisotropic morphological Gaussian kernels, producing edge strength and direction maps.

\textbf{Step 5: Edge fusion.} The edge maps are combined to obtain a refined and continuous edge representation.

\textbf{Step 6: Post-processing.}
\begin{itemize}
    \item Noise suppression,
    \item Edge thinning,
    \item Removal of spurious responses.
\end{itemize}

\textbf{Step 7: Line extraction and measurement.} Detected edges are further processed to extract separation lines between rubber strips and compute distances between them.

\subsubsection{Output}

The output of the proposed pipeline includes:
\begin{itemize}
    \item binary edge map,
    \item continuous edge contours,
    \item detected separation lines,
    \item pixel coordinates of edge boundaries.
\end{itemize}

\subsubsection{Evaluation metrics}

The performance of the method is evaluated using standard edge detection metrics:
\begin{itemize}
    \item Precision and Recall,
    \item F1-score,
    \item Pratt’s Figure of Merit (FOM),
    \item Precision-Recall (PR) curves.
\end{itemize}

These metrics assess both edge localization accuracy and robustness to noise.

\subsubsection{Applicability to the Michelin Challenge}

The selected method is particularly suitable for the Michelin Challenge 2 because it provides robust edge detection in noisy industrial environments without requiring training data or GPU resources. Its multi-scale and anisotropic design allows accurate detection of thin separation lines between rubber strips, even under low contrast and irregular surface conditions. This makes it an efficient and deployable solution for automated inspection systems, where reliability, interpretability, and computational efficiency are critical.

\subsection{Adaptive Classical Vision Pipeline for Cut-Edge Detection and Metric Estimation}

\subsubsection{Description and Motivation}

The selected method is an adaptive classical computer vision pipeline for industrial edge detection and metric estimation, built upon the improved Canny edge detection framework proposed by Liu et al.~\cite{Liu2024b} as its primary methodological foundation. This approach was selected because it directly addresses the core challenge of the Michelin Challenge 2, namely the detection of thin, approximately linear cut edges on a rubber strip, without requiring any training data or GPU-based infrastructure. As a result, it can be deployed immediately in industrial environments with minimal setup.

The pipeline is extended with techniques drawn from additional reviewed works: adaptive image segmentation and contour smoothing from González et al.~\cite{Gonzalez2024}, perspective correction and pixel-to-millimetre calibration from Ding et al.~\cite{Ding2025Homography} and Zhang et al.~\cite{Zhang2024Calibration}, and a grid-based region of interest strategy from Liu et al.~\cite{Liu2024a}. This combination enables the method to cover the complete processing chain required by the challenge, from raw smartphone image acquisition to metric estimation of cut-edge positions in millimetres.

\subsubsection{Computer Vision Pipeline}

The proposed pipeline consists of five sequential stages:

\paragraph{Step 1 --- Image preprocessing}
The input is a raw RGB image captured with a smartphone under varying lighting conditions and camera perspectives. The image is first converted to grayscale, reducing computational complexity. Following González et al.~\cite{Gonzalez2024}, Contrast Limited Adaptive Histogram Equalization (CLAHE) is applied instead of global histogram equalization. This technique operates on local regions, preventing over-amplification of noise in uniform areas, which is crucial in uncontrolled industrial lighting conditions. Finally, a Gaussian blur with a $5 \times 5$ kernel is applied to suppress high-frequency noise prior to edge detection, following standard Canny-based practices~\cite{Liu2024b}.

\paragraph{Step 2 --- QR code detection and perspective rectification}
The three QR codes present in the image are detected using the \texttt{pyzbar} library, which relies on classical binarization and pattern matching. The pixel coordinates of the QR code centres are used to compute a homography matrix via \texttt{cv2.findHomography} with RANSAC-based outlier rejection, following the approach of Ding et al.~\cite{Ding2025Homography}. Given that the QR codes form a right triangle with known dimensions (800 mm), the homography maps image coordinates to a metric coordinate system where $(0,0)$ corresponds to the upper-left QR code.

The image is then rectified using \texttt{cv2.warpPerspective}, correcting perspective distortion caused by camera positioning. Additionally, lens distortion is corrected using the single-image calibration method proposed by Zhang et al.~\cite{Zhang2024Calibration}, which estimates radial and tangential distortion parameters.

\paragraph{Step 3 --- Adaptive Canny edge detection}
Edge detection is performed using an improved version of the Canny operator based on Liu et al.~\cite{Liu2024b}. Instead of manually selecting thresholds, Otsu’s method is used to compute an optimal threshold $t_{\text{Otsu}}$ from the grayscale histogram. The Canny thresholds are then defined as:
\[
t_{\text{low}} = 0.5 \cdot t_{\text{Otsu}}, \quad t_{\text{high}} = t_{\text{Otsu}}
\]
This adaptive strategy ensures robustness to variations in illumination and contrast. Internally, the Canny algorithm computes gradients using Sobel operators, applies non-maximum suppression to obtain thin edges, and performs hysteresis thresholding to retain strong and connected edge responses.

\paragraph{Step 4 --- Region of interest extraction and edge filtering}
The binary edge map produced by the Canny detector contains edges from multiple structures. To focus on relevant regions, a grid-based region of interest (ROI) strategy is applied, inspired by Liu et al.~\cite{Liu2024a}, where regions with low edge density are discarded.

Within the selected regions, the Probabilistic Hough Transform (\texttt{cv2.HoughLinesP}) is used to extract line segments. A geometric filtering stage is then applied: segments are filtered based on minimum length, orientation, and spatial position relative to the QR-based coordinate system. Collinear segments within a proximity threshold are merged into a single representative line, following a contour smoothing strategy similar to González et al.~\cite{Gonzalez2024}.

\paragraph{Step 5 --- Metric estimation}
Once the cut-edge lines are identified, 10 equally spaced sampling points are extracted along each detected edge. Since the image has been rectified into a metric coordinate system, pixel coordinates correspond directly to physical positions in millimetres. The distance from each sampling point to the lower border of the table is computed using Euclidean distance in this metric space.

\subsubsection{Output}

The pipeline produces the following outputs:
\begin{itemize}
    \item Pixel coordinates of the detected cut-edge boundaries.
    \item Metric positions of the edges in the QR-based coordinate system, where $(0,0)$ corresponds to the upper-left QR code.
    \item Ten distance measurements per edge, corresponding to sampling points, expressed in millimetres.
\end{itemize}

\subsubsection{Evaluation Metrics}

The performance of the pipeline is evaluated using several metrics. For edge detection accuracy, the Root Mean Square Error (RMSE) between detected edges and ground truth annotations is used, following González et al.~\cite{Gonzalez2024}. For segmentation quality, Intersection over Union (IoU) is computed between predicted and ground truth masks. Precision and Recall are also reported to assess the trade-off between false positives and false negatives.

Processing speed is measured in seconds per image on a standard CPU, evaluating the practical feasibility of deploying the pipeline in real industrial environments.

\subsubsection{Applicability to the Michelin Challenge}

The proposed pipeline is particularly suitable for the Michelin Challenge 2 as it relies entirely on classical image processing techniques implemented in OpenCV and standard Python libraries. It enables QR-based spatial calibration, perspective correction, robust edge detection under varying conditions, and accurate metric estimation without requiring training data or specialised hardware.

The adaptive thresholding strategy ensures robustness to variations in lighting and camera positioning, while the homography-based transformation allows precise millimetre-level measurements. This makes the method highly practical for deployment in real factory environments using only a smartphone camera.
% TODO: \subsection{Traditional Method (to be completed by group member)}

% --------------------
\section{Datasets}
% --------------------
\label{sec:datasets}
% At least two publicly available datasets per student, in paragraph format.
% Include name, date, link, and description for each dataset.

Several publicly available datasets are commonly used to evaluate object detection and 
segmentation models such as Grounding DINO and the Segment Anything Model. In 
addition, this project makes use of a proprietary dataset provided by Michelin 
Aranda de Duero as part of the challenge.

The \textbf{COCO (Common Objects in Context)} dataset was introduced in 2014 by 
Microsoft~\cite{LinCOCO2014} and is available at \url{https://cocodataset.org}. 
COCO contains more than 330,000 images, with over 200,000 labeled images and more 
than 1.5 million object instances annotated across 80 object categories, covering 
common objects in everyday scenes such as people, vehicles, animals, and household 
items. Each image provides detailed annotations including bounding boxes, instance 
segmentation masks, and keypoints, making the dataset suitable for evaluating object 
detection, instance segmentation, and keypoint detection methods. Due to its diversity 
and scale, COCO is the primary benchmark for evaluating modern detection models such 
as Grounding DINO and YOLOv8.Two samples from the dataset with their segmentation are shown in Figure~\ref{fig:coco_samples}\\

\begin{figure}[htb!]
    \centering
    \begin{subfigure}[b]{0.48\linewidth}
        \includegraphics[width=\linewidth, height=5cm, keepaspectratio]
        {figures/coco_times_original.jpg}
    \end{subfigure}
    \hfill
    \begin{subfigure}[b]{0.48\linewidth}
        \includegraphics[width=\linewidth, height=5cm, keepaspectratio]
        {figures/coco_times_segmented.jpg}
    \end{subfigure}
    \vspace{0.3cm}
    \begin{subfigure}[b]{0.48\linewidth}
        \includegraphics[width=\linewidth, height=5cm, keepaspectratio]
        {figures/coco_dog_original.jpg}
    \end{subfigure}
    \hfill
    \begin{subfigure}[b]{0.48\linewidth}
        \includegraphics[width=\linewidth, height=5cm, keepaspectratio]
        {figures/coco_dog_segmented.jpg}
    \end{subfigure}
    \caption{Sample images from the COCO dataset~\cite{LinCOCO2014}. 
    Left column: original images. Right column: instance segmentation 
    annotations showing object masks across multiple categories.}
    \label{fig:coco_samples}
\end{figure}

Another relevant dataset is \textbf{NEU Surface Defect Database} was introduced in 2013 by Northeastern 
University~\cite{NeuDataset} and is available at 
\url{https://ieee-dataport.org/documents/neu-det}. 
The dataset contains 1,800 grayscale images of hot-rolled steel strip surfaces, 
divided into six categories of surface defects: rolled-in scale, patches, crazing, 
pitted surface, inclusion, and scratches. Each category contains 300 samples, and 
the images have a resolution of $200 \times 200$ pixels. The dataset provides both 
image-level labels for classification tasks and bounding box annotations for object 
detection tasks, making it particularly suitable for evaluating defect detection 
methods in industrial inspection scenarios. Due to its focus on surface defect 
localization in manufacturing environments, NEU is directly relevant to the 
Michelin Challenge, where precise detection and localization of structural 
anomalies in rubber strips is required. Some images of the dataset are shown in Figure~\ref{fig:neu_samples}\\

\begin{figure}[htb!]
    \centering
    \begin{subfigure}[b]{0.48\linewidth}
        \includegraphics[width=\linewidth, height=5cm, keepaspectratio]
        {figures/neu_sample1.png}
    \end{subfigure}
    \hfill
    \begin{subfigure}[b]{0.48\linewidth}
        \includegraphics[width=\linewidth, height=5cm, keepaspectratio]
        {figures/neu_sample2.png}
    \end{subfigure}
    \caption{Sample images from the NEU Surface Defect Database~\cite{NeuDataset}, 
    showing bounding box annotations for surface defect detection on hot-rolled 
    steel strips.}
    \label{fig:neu_samples}
\end{figure}

The \textbf{LVIS (Large Vocabulary Instance Segmentation)} dataset was introduced 
in 2019 by Gupta et al. at Facebook AI Research and is available at its official 
website\footnote{\url{https://www.lvisdataset.org}}. LVIS contains over 164,000 
images sourced from the COCO image collection, annotated with high-quality instance 
segmentation masks across 1,203 object categories, making it one of the most 
diverse and fine-grained instance segmentation benchmarks available. The dataset 
is divided into three frequency-based splits: frequent categories (more than 100 
training samples), common categories (between 10 and 100 samples), and rare 
categories (fewer than 10 samples), reflecting the long-tail distribution of 
object categories in the real world. Annotations were collected using a 
federated annotation protocol that assigns only a subset of categories to each 
image, resulting in over 2 million high-quality segmentation masks. LVIS is used 
to evaluate EfficientSAM \cite{Xiong2023EfficientSAM} on zero-shot instance 
segmentation, where EfficientSAM-S achieves 42.3 AP on LVIS compared to 44.7 AP 
for the original SAM, while using $20\times$ fewer parameters. Two sample images from the dataset are shown in Figure~\ref{fig:LVIS_samples}.\\

%% FALTA AÑADIR LAS IMAGENES

The \textbf{SA-1B} dataset was introduced in 2023 by Kirillov et al. at Meta FAIR 
and is available at its official 
website\footnote{\url{https://ai.meta.com/datasets/segment-anything/}}. SA-1B is 
the largest image segmentation dataset to date, comprising over 11 million 
high-resolution images with more than 1.1 billion automatically generated segmentation 
masks. The masks cover objects of all categories without semantic constraints, 
including whole objects and their parts and subparts, making it a general-purpose 
dataset for training and evaluating promptable segmentation models. It is used to 
pretrain and evaluate both SAM 2 \cite{Ravi2024SAM2} and EfficientSAM 
\cite{Xiong2023EfficientSAM}, with the latter achieving 44.4 AP on COCO zero-shot 
instance segmentation after fine-tuning on SA-1B.\\

% FALTA AÑADIR LAS IMAGENES

The BSDS500 (Berkeley Segmentation Dataset and Benchmark) was introduced by the University of California, Berkeley, and is available at \url{https://www2.eecs.berkeley.edu/Research/Projects/CS/vision/bsds/} The dataset contains a total of 500 natural images, divided into 200 training, 100 validation, and 200 testing images. Each image is annotated by multiple human subjects, providing ground truth boundaries that reflect perceptual edge information \cite{Martin2001BSDS,Arbelaez2011ContourDetection}.\\

Another relevant dataset is PASCAL VOC 2007 (Visual Object Classes), introduced as part of the PASCAL challenge and available at \url{https://www.robots.ox.ac.uk/~vgg/projects/pascal/VOC/voc2007/} The dataset contains 9,963 images with approximately 24,640 annotated object instances across 20 object categories, including people, vehicles, animals, and everyday objects \cite{Everingham2010PASCALVOC}.\\

Lastly, we have \textbf{Michelin Aranda Dataset}, which is a propietary image collection provided by
Michelin Aranda. The dataset consists of images 
captured in laboratory conditions simulating real factory scenarios, using a flat table 
as a surface. Each image contains rubber strips placed on the table alongside three QR 
codes arranged in a right triangle with legs of 800 mm, used as calibration markers 
for pixel-to-metric conversion. The images were captured with a smartphone under 
variable camera positions and angles, reflecting the real acquisition conditions on 
the production floor. The dataset is specifically designed for the tasks of QR 
detection, rubber strip segmentation, cut edge localization, and metric distance 
estimation. Two sample images from this dataset are shown in 
Figure~\ref{fig:michelin_samples}.

\begin{figure}[htb!]
    \centering
    \begin{subfigure}[b]{0.48\linewidth}
        \includegraphics[width=\linewidth, height=6cm, keepaspectratio]{figures/michelin1.jpg}
    \end{subfigure}
    \hfill
    \begin{subfigure}[b]{0.48\linewidth}
        \includegraphics[width=\linewidth, height=6cm, keepaspectratio]{figures/michelin2.jpg}
    \end{subfigure}
    \caption{Sample images from the Michelin Aranda Dataset. QR codes are visible 
    in the corners of the surface and serve as spatial calibration markers.}
    \label{fig:michelin_samples}
\end{figure}

% --------------------
% --------------------
\section{Evaluation Metrics}
% --------------------
\label{sec:metrics}

The performance of the proposed system is evaluated using standard metrics commonly 
adopted in object detection, image segmentation, and edge detection tasks. The 
following metrics are used across the four pipelines proposed in this work.

\paragraph{Precision.}
Precision measures the proportion of correctly predicted detections among all 
predicted detections, reflecting the model's ability to avoid false positives.

\[
\text{Precision} = \frac{TP}{TP + FP}
\]

\paragraph{Recall.}
Recall measures the ability of the model to detect all relevant objects present 
in the image, reflecting the model's capacity to minimise false negatives.

\[
\text{Recall} = \frac{TP}{TP + FN}
\]

\paragraph{F1-score.}
The F1-score provides a balanced measure between Precision and Recall, combining 
both into a single metric. It is particularly useful when there is an uneven class 
distribution.

\[
F1 = \frac{2 \cdot \text{Precision} \cdot \text{Recall}}{\text{Precision} + \text{Recall}}
\]

\paragraph{Intersection over Union (IoU).}
Intersection over Union evaluates the overlap between the predicted region and the 
ground truth annotation, defined as the ratio between the intersection area and 
the union area of both regions.

\[
\text{IoU} = \frac{\text{Area of Overlap}}{\text{Area of Union}}
\]

A prediction is typically considered correct if the IoU exceeds a predefined 
threshold (e.g., 0.5), although stricter evaluations may use multiple thresholds. 
IoU is used in both object detection and segmentation tasks to assess the quality 
of predicted masks and bounding boxes.

\paragraph{Average Precision (AP).}
Average Precision summarises the precision-recall curve for a given class by 
computing the area under the curve, providing a single value representing model 
performance across different recall levels.

\[
\text{AP} = \int_0^1 \text{Precision}(\text{Recall}) \, d(\text{Recall})
\]

\paragraph{Mean Average Precision (mAP).}
Mean Average Precision is computed as the mean of the Average Precision values 
across all object classes, and is one of the most widely used metrics in object 
detection benchmarks.

\[
\text{mAP} = \frac{1}{N} \sum_{i=1}^{N} AP_i
\]

where $N$ is the number of object classes and $AP_i$ is the Average Precision 
for class $i$. In this work, mAP is used to evaluate the detection of QR codes 
and rubber strip cut edges by Grounding DINO and YOLOv9.

\paragraph{Precision-Recall (PR) Curve.}
The PR curve represents the trade-off between Precision and Recall at different 
classification thresholds. It is widely used in edge detection benchmarks such 
as BSDS500 to provide a comprehensive view of detector performance across 
operating points.

\paragraph{Pratt's Figure of Merit (FOM).}
Pratt's FOM evaluates edge localization accuracy by penalising both missed 
detections and spatial deviations between detected edges and ground truth edges.

\[
\text{FOM} = \frac{1}{\max(N_d, N_g)} \sum_{i=1}^{N_d} \frac{1}{1 + \alpha d_i^2}
\]

where $N_d$ is the number of detected edge pixels, $N_g$ is the number of ground 
truth edge pixels, $d_i$ is the distance between a detected edge pixel and the 
nearest ground truth edge pixel, and $\alpha$ is a scaling constant. FOM is used 
to evaluate the MAMDD classical edge detector.

\paragraph{Root Mean Square Error (RMSE).}
RMSE measures the average magnitude of errors between predicted metric positions 
and ground truth measurements, penalising large deviations more heavily than 
small ones.

\[
\text{RMSE} = \sqrt{\frac{1}{n} \sum_{i=1}^{n} (\hat{y}_i - y_i)^2}
\]

where $\hat{y}_i$ is the predicted metric position of the $i$-th point and $y_i$ 
is the corresponding ground truth value. RMSE is the primary metric for evaluating 
the millimetre-level accuracy of the adaptive Canny pipeline.

\paragraph{Mean Absolute Error (MAE).}
MAE measures the average absolute difference between predicted and ground truth 
metric positions. Unlike RMSE, it treats all errors equally regardless of their 
magnitude, providing a more interpretable measure of average deviation in 
millimetres.

\[
\text{MAE} = \frac{1}{n} \sum_{i=1}^{n} |\hat{y}_i - y_i|
\]

\paragraph{Distance Error.}
The Distance Error evaluates geometric accuracy by measuring the absolute 
difference between the predicted and ground truth distance between key elements, 
such as cutting edges and reference QR markers.

\[
\text{Distance Error} = | d_{\text{pred}} - d_{\text{gt}} |
\]

where $d_{\text{pred}}$ is the estimated distance and $d_{\text{gt}}$ is the 
ground truth measurement in millimetres. This metric is particularly important 
for the Michelin Challenge, where precise spatial measurements are required for 
industrial automation.\\

In summary, Precision, Recall, F1-score, and PR curves are used to evaluate the 
quality of edge detection across all pipelines. IoU and mAP assess the accuracy 
of segmentation masks and object detections respectively. FOM provides a 
comprehensive assessment of edge continuity and localization accuracy for the 
classical methods. RMSE, MAE, and Distance Error directly evaluate the precision 
of millimetre-level geometric measurements, which are the primary output required 
by the Michelin Challenge.

% --------------------
% --------------------
% --------------------
\section{Python Implementations}
% --------------------
\label{sec:python}

This section presents the available Python implementations of the methods reviewed 
in Section~\ref{sec:literatureReview}, including their repositories, release dates, 
main advantages, and limitations, with a focus on their practical applicability to 
the Michelin Challenge 2.

% --------------------
\subsection{Deep Learning Methods}
% --------------------

\textbf{YOLOv8}~\cite{Jocher2023YOLOv8} is available at 
\url{https://github.com/ultralytics/ultralytics}. The official implementation was 
released on January 10, 2023, by Ultralytics as a unified Python package available 
via pip, built on PyTorch with both a Python API and a command-line interface. It 
supports multiple model sizes (n, s, m, l, x) and export to ONNX, TensorRT, and 
CoreML formats. Its main advantages are ease of use, real-time inference capability, 
and support for advanced data augmentation strategies such as mosaic and mixup. Its 
main limitations are a closed vocabulary that restricts detection to predefined 
categories, and significantly higher computational requirements for larger variants.\\

\textbf{RT-DETRv4}~\cite{LiaoRTDETRv4} is available at 
\url{https://github.com/RT-DETRs/RT-DETRv4}. The repository was created on 
October 30, 2025, with full code, configurations, and checkpoints released on 
November 17, 2025. The implementation is provided in PyTorch and includes training, 
evaluation, and deployment scripts for four model sizes (S, M, L, X) using HGNetv2 
as backbone. Its main advantages are high detection accuracy through transformer-based 
global attention and an end-to-end pipeline that avoids Non-Maximum Suppression. Its 
main limitations are higher computational requirements than lightweight CNN detectors 
and greater implementation complexity.\\

\textbf{AnoDDPM}~\cite{WyattAnoDDPM2022} is available at 
\url{https://github.com/Julian-Wyatt/AnoDDPM}. The method was introduced at CVPR 
Workshops 2022 and is written in PyTorch, providing scripts for training on normal 
data and running anomaly detection at inference time. Note that the repository was 
archived by the author in October 2025 and is no longer actively maintained. Its 
main advantages are the ability to learn exclusively from defect-free data and stable 
training compared to GAN-based approaches. Its main limitations are high computational 
cost due to the iterative denoising process and the requirement for careful tuning of 
noise schedules.\\

\textbf{Segformer++}~\cite{KienzleSegformer2024} is available at 
\url{https://github.com/KieDani/SegformerPlusPlus}. It was introduced in May 2024 
and is built on PyTorch with integration with the MMSegmentation and MMPose 
libraries. A key practical advantage is that the token merging strategy can be 
applied at inference time using the original Segformer pre-trained weights without 
any retraining. Its main limitations are that performance gains are less pronounced 
at lower resolutions, and the method is specifically adapted to the Segformer 
architecture.\\

\textbf{PEM}~\cite{CavagneroPEM2024} is available at 
\url{https://github.com/NiccoloCavagnero/PEM}. It was introduced on February 29, 
2024, and is provided in PyTorch with support for both semantic and panoptic 
segmentation under a single unified codebase. Its main advantages are significantly 
faster inference than comparable models such as Mask2Former and a unified framework 
that avoids task-specific architectures. Its main limitations are that performance 
falls slightly below the most accurate but slower models on some benchmarks, and 
computational requirements remain higher than lightweight CNN-only architectures.\\

\textbf{Relation-DETR}~\cite{HouRelationDETR2024} is available at 
\url{https://github.com/xiuqhou/Relation-DETR}. It was introduced on July 16, 2024, 
and is provided in PyTorch with full training and evaluation scripts for COCO 2017. 
The position relation encoder is designed as a plug-in-and-play module compatible 
with any DETR-based architecture. Its main advantages are significantly faster 
convergence and state-of-the-art detection performance. Its main limitations are 
increased memory requirements at training time and limited validation on 
domain-specific datasets.\\

\textbf{ViLU-Net}~\cite{HeidariViLUNet2025} is available at 
\url{https://github.com/moeinheidari7829/Retroperitoneal_Tumour_Segmentation_SPIE2025}. 
It was introduced on February 1, 2025, and is built on top of the nnU-Net framework 
in PyTorch. Its main advantages are state-of-the-art performance on medical 
segmentation benchmarks and efficient long-range dependency modeling via xLSTM 
blocks. Its main limitations are the small in-house dataset used for evaluation, 
high GPU memory requirements (NVIDIA V100, 32 GB), and limited validation beyond 
CT imaging.\\

\textbf{Grounded-SAM}~\cite{RenGroundedSAM2024} is available at 
\url{https://github.com/IDEA-Research/Grounded-Segment-Anything}. It was introduced 
on January 25, 2024, and is actively maintained. The implementation follows a 
training-free pipeline combining Grounding DINO with SAM, and includes ready-to-use 
scripts for multiple pipelines including integration with HQ-SAM, MobileSAM, and 
EfficientSAM. Its main advantages are open-vocabulary detection and segmentation 
without retraining and strong zero-shot performance on the SegInW benchmark. Its 
main limitations are inference speed bottlenecks for real-time applications and 
dependence on the quality of Grounding DINO detections.\\

\textbf{YOLO-World}~\cite{ChengYOLOWorld2024} is available at 
\url{https://github.com/AILab-CVC/YOLO-World}. It was introduced on January 30, 
2024, and is built on top of the MMYOLO and MMDetection toolboxes in PyTorch. It 
provides three model variants (S, M, L) and supports pre-training on large-scale 
datasets as well as fine-tuning for downstream tasks. Its main advantages are 
real-time open-vocabulary detection at 52 FPS and a lightweight architecture with 
as few as 13M parameters. Its main limitations are lower zero-shot performance than 
heavier open-vocabulary detectors and high pre-training computational requirements.\\

\textbf{Grounding DINO 1.5}~\cite{RenGroundingDINO15} is available at 
\url{https://github.com/IDEA-Research/Grounding-DINO-1.5-API}. It was introduced 
on May 16, 2024, and is accessible via API rather than as a fully open-source 
codebase, with two model variants: the Pro model using a ViT-L backbone for maximum 
accuracy, and the Edge model using EfficientViT-L1 for real-time inference. Its main 
advantages are state-of-the-art zero-shot detection performance and support for 
diverse input modalities. Its main limitations are the closed API restricting 
reproducibility and customization, and the high computational requirements of the 
Pro model.\\

\textbf{SAM 2}~\cite{Ravi2024SAM2} is available at 
\url{https://github.com/facebookresearch/sam2}. The official implementation was 
released on August 1, 2024, by Meta FAIR as an open-source PyTorch package under 
an Apache 2.0 license. It supports both image and video segmentation through a 
unified promptable interface accepting points, bounding boxes, and masks, with 
multiple model sizes available (Hiera-T, S, B+, L). Its main advantages are strong 
zero-shot segmentation performance without task-specific training and a streaming 
memory architecture enabling video processing. Its main limitations are high GPU 
memory requirements for larger variants and class-agnostic mask outputs that require 
additional classification logic for specific downstream tasks.\\

\textbf{EfficientSAM}~\cite{Xiong2023EfficientSAM} is available at 
\url{https://github.com/yformer/EfficientSAM}. The implementation was released in 
December 2023 by Meta AI Research in PyTorch, providing two model variants: 
EfficientSAM-Ti (10M parameters) and EfficientSAM-S (25M parameters), both 
pretrained via SAMI on SA-1B. Its main advantages are a $20\times$ reduction in 
parameter count and inference time compared to SAM ViT-H while maintaining 
competitive segmentation quality. Its main limitations are reduced documentation 
compared to SAM 2, the absence of video segmentation support, and lower performance 
on fine-grained segmentation tasks.\\

\textbf{YOLOv9}~\cite{Wang2024YOLOv9} is available at 
\url{https://github.com/WongKinYiu/yolov9}. The implementation was released on 
February 29, 2024, in PyTorch with full training, evaluation, and inference scripts 
for MS COCO, providing multiple model sizes (S, M, C, E) and supporting both object 
detection and instance segmentation. Its main advantages are state-of-the-art 
accuracy among train-from-scratch detectors and well-documented training pipelines 
that facilitate fine-tuning on custom datasets. Its main limitations are that the 
PGI auxiliary branch adds training complexity, and the repository has seen reduced 
maintenance activity compared to the Ultralytics ecosystem.\\

\textbf{YOLOv8-HDSA}~\cite{Shi2025CarbonBlock} is based on the Ultralytics YOLOv8 
framework, available at \url{https://github.com/ultralytics/ultralytics}. The base 
framework was released in January 2023 and is actively maintained, with the 
HDSA-specific modifications introduced in the paper published in March 2025. Its 
main advantages are an exceptionally mature ecosystem with extensive documentation 
and easy fine-tuning on custom datasets via a single configuration file. Its main 
limitation is that the HDSA-specific modifications (SRFF module, convolutional 
self-attention head, Focaler-IoU loss) are not released separately and must be 
reimplemented from the paper description.\\

\textbf{AnomalyDINO}~\cite{Damm2024AnomalyDINO} is available at 
\url{https://github.com/SimonDamm/AnomalyDINO}. The implementation was released in 
May 2024 and accepted at WACV 2025, leveraging DINOv2 pretrained weights from 
\texttt{torch.hub} without requiring additional training data beyond a few reference 
normal images. Its main advantages are a completely training-free inference pipeline 
and state-of-the-art one-shot performance achieving $96.6\%$ AUROC on MVTec-AD. Its 
main limitations are performance degradation when fewer than 3--5 reference images 
are available and sensitivity to domain shift between reference and test images.\\

\textbf{CLIP-DINOv2}~\cite{Jiang2025CLIPDINO} does not have a dedicated public 
repository. The paper, published in December 2025, provides full implementation 
details including architecture specifications, hyperparameters, and training 
procedures. The method relies on publicly available pretrained weights for CLIP 
ViT-L/14@336px from \url{https://github.com/openai/CLIP} and DINOv2 ViT-L from 
\url{https://github.com/facebookresearch/dinov2}. Its main advantages are 
state-of-the-art zero-shot anomaly detection performance across seven industrial 
benchmarks. Its main limitation is the absence of a public codebase, requiring full 
reimplementation from the paper to reproduce results.\\

\textbf{SAM2-UNet}~\cite{Xiong2026SAM2UNet} is available at 
\url{https://github.com/WZH0120/SAM2-UNet}. The peer-reviewed version was published 
in Visual Intelligence in January 2026, implemented in PyTorch and requiring SAM 2 
pretrained Hiera weights from \url{https://github.com/facebookresearch/sam2}. Its 
main advantages are parameter-efficient fine-tuning via lightweight adapters that 
keep the Hiera encoder frozen, enabling training on a single 24 GB GPU and strong 
generalization across diverse segmentation tasks. Its main limitation is that the 
current implementation operates at $352\times352$ resolution, which may reduce 
performance on high-resolution images with fine edge details.\\

\textbf{Semi-supervised defect segmentation}~\cite{Shi2024SemiSupervised} does not 
have a dedicated public repository. The paper was published in Scientific Reports in 
September 2024, with code available upon request from the corresponding author via 
\url{https://doi.org/10.1038/s41598-024-72579-6}. The implementation is based on 
PyTorch with a modified DeepLabv3Plus backbone using a simplified MobileNetv2 
encoder. Its main advantages are a lightweight architecture achieving 64.59 FPS 
while outperforming state-of-the-art semi-supervised methods and a dynamic threshold 
strategy that improves pseudo-label quality over training. Its main limitation is 
the lack of a public repository, which hinders reproducibility.\\

\textbf{Camera Calibration Survey}~\cite{Liao2025CameraCalibration} maintains a 
curated repository at 
\url{https://github.com/KangLiao929/Awesome-Deep-Camera-Calibration}, last revised 
in February 2025. The repository is not a single codebase but a structured 
collection of links to calibration method implementations organized by category, 
alongside a unified holistic calibration dataset. Its main advantage is providing a 
comprehensive reference for selecting and comparing calibration methods relevant to 
the QR-based spatial calibration required in the Michelin Challenge. Its main 
limitation is that it does not provide ready-to-use calibration code and requires 
selecting and integrating individual method repositories.\\

\textbf{MVTec AD 2}~\cite{HecklerKram2025MVTecAD2} provides the dataset and 
evaluation scripts at \url{https://www.mvtec.com/company/research/datasets/mvtec-ad-2}. 
Released in March 2025 by MVTec Software GmbH, it does not provide a segmentation 
or detection model but offers a standardized evaluation framework compatible with 
existing anomaly detection pipelines. Its main advantage is providing a challenging 
and realistic benchmark covering scenarios not addressed by existing datasets. Its 
main limitation is that no baseline model is included, requiring users to integrate 
their own methods for evaluation.\\

The deep learning methods reviewed span a wide spectrum of approaches. 
Foundation segmentation models --- SAM 2, EfficientSAM, and SAM2-UNet --- share a 
common architectural lineage based on the Hiera encoder and promptable segmentation 
paradigm, differing mainly in computational efficiency and downstream adaptability. 
Object detection methods --- YOLOv8, RT-DETRv4, Relation-DETR, YOLO-World, 
Grounding DINO 1.5, and Grounded-SAM --- range from efficient real-time CNN-based 
detectors to transformer-based open-vocabulary models. Industrial defect methods 
--- YOLOv9, YOLOv8-HDSA, AnomalyDINO, CLIP-DINOv2, and the semi-supervised 
approach --- differ fundamentally in their supervision requirements, from fully 
supervised fine-tuning to completely training-free inference. From the perspective 
of the Michelin Challenge, where the available dataset is small and annotations are 
limited, the most practical deep learning approaches are AnomalyDINO for its 
training-free nature and the Grounded-SAM and YOLOv9~+~EfficientSAM pipelines for 
their publicly available, well-documented, and PyTorch-compatible implementations.\\

% --------------------
\subsection{Classical Methods}
% --------------------

\textbf{Canny Edge Detector}~\cite{Canny1986} is available at 
\url{https://github.com/opencv/opencv}. The implementation is part of the OpenCV 
library, originally released in 2000 and continuously maintained, accessible in 
Python through the \texttt{cv2.Canny()} function. Its main advantages are 
simplicity, efficiency, and real-time performance, making it a standard baseline 
for edge detection. Its main limitations are sensitivity to noise and fragmented 
edge responses in complex textures such as rubber surfaces.\\

\textbf{Line Segment Detector (LSD)}~\cite{Grompone2010LSD} is available at 
\url{https://github.com/opencv/opencv}. The implementation is included in OpenCV 
via the \texttt{cv2.createLineSegmentDetector()} function. LSD directly detects 
line segments instead of raw edges, which is particularly useful for identifying 
separation lines. Its main advantages are high speed, no need for parameter tuning, 
and direct extraction of geometric structures. Its main limitations are sensitivity 
to weak edges and reduced performance in very low-contrast scenarios.\\

\textbf{Structured Edge Detection}~\cite{Dollar2013StructuredEdges} is available 
at \url{https://github.com/opencv/opencv_extra}, implemented via the 
\texttt{cv2.ximgproc.createStructuredEdgeDetection()} function using pre-trained 
structured forest models. The method provides improved edge continuity and 
robustness compared to classical detectors. Its main advantages are better edge 
quality and robustness to noise. Its main limitations are the need for pre-trained 
models and slightly higher computational cost compared to Canny.\\

\textbf{Watershed Segmentation}~\cite{Vincent1991Watershed} is available at 
\url{https://github.com/opencv/opencv}, implemented via the 
\texttt{cv2.watershed()} function. It is widely used for region-based segmentation 
and object separation. Its main advantages are the ability to separate touching 
regions and its interpretability. Its main limitations are over-segmentation and 
sensitivity to marker initialization.\\

\textbf{Frangi Filter}~\cite{Frangi1998Vesselness} is available at 
\url{https://github.com/scikit-image/scikit-image}, implemented via the 
\texttt{skimage.filters.frangi} function. The method is designed to enhance 
line-like structures, making it useful for detecting thin separations. Its main 
advantages are strong performance on thin structures and robustness to noise. Its 
main limitations are sensitivity to parameter selection and increased computational 
cost.\\

\textbf{OpenCV}~\cite{opencv_library} is available at 
\url{https://github.com/opencv/opencv}. OpenCV is an open-source computer vision 
library first released in 2000 and actively maintained. It provides native Python 
bindings and includes all the classical image processing operations required by the 
adaptive Canny pipeline, such as Gaussian blurring, Otsu thresholding, Canny edge 
detection, the Probabilistic Hough Transform, homography computation via 
\texttt{cv2.findHomography}, perspective rectification via 
\texttt{cv2.warpPerspective}, and lens distortion correction via 
\texttt{cv2.undistort}. Its main advantages are maturity, extensive documentation, 
cross-platform support, and independence from GPU hardware. Its main limitation is 
the absence of high-level abstractions for pipeline orchestration, requiring manual 
integration of each processing stage.\\

\textbf{pyzbar} is available at 
\url{https://github.com/NaturalHistoryMuseum/pyzbar}, released in 2016 and actively 
maintained. This library wraps the ZBar framework and provides a simple interface 
for detecting and decoding QR codes from images and video streams. Its main 
advantage is simplicity and reliability without requiring any training data. Its 
main limitation is reduced robustness under strong perspective distortion or partial 
occlusions.\\

\textbf{Grid-HOG-PCA-LightGBM}~\cite{Liu2024a} is implemented using 
\texttt{scikit-image} for HOG feature extraction 
(\url{https://github.com/scikit-image/scikit-image}), \texttt{scikit-learn} for 
PCA dimensionality reduction 
(\url{https://github.com/scikit-learn/scikit-learn}), and LightGBM 
(\url{https://github.com/microsoft/LightGBM}) for classification. These libraries 
are widely used, well-documented, and installable via \texttt{pip} without GPU 
requirements. Its main limitation is that HOG feature extraction is sensitive to 
parameter choices such as cell size and block normalization, requiring tuning for 
different datasets.\\

\textbf{Improved Canny}~\cite{Liu2024b} can be directly implemented using OpenCV 
functions \texttt{cv2.Canny} and \texttt{cv2.threshold} with the 
\texttt{THRESH\_OTSU} flag, requiring no additional repository beyond the standard 
OpenCV distribution. Its main advantage is simplicity and immediate availability. 
Its main limitation is that Otsu-based thresholding assumes a bimodal intensity 
distribution, which may not always hold in complex industrial scenes.\\

\textbf{Homography-based calibration}~\cite{Ding2025Homography,Zhang2024Calibration} 
is implemented using \texttt{cv2.findHomography} and \texttt{cv2.calibrateCamera}, 
both included in OpenCV. A reference implementation is available at 
\url{https://docs.opencv.org/4.x/dc/dbb/tutorial_py_calibration.html}. Its main 
advantage is that the full calibration and rectification pipeline requires no 
external dependencies. Its main limitation is that calibration accuracy depends on 
the quality of detected keypoints and may degrade with low-resolution or poorly 
captured images.\\

The classical methods reviewed can be grouped into two categories: edge detection 
and structure extraction. Edge detection methods --- Canny, Structured Edge 
Detection, and the improved Canny pipeline --- focus on identifying intensity 
discontinuities and boundaries, differing mainly in their robustness to noise and 
the degree of parameter automation. Structure extraction methods --- LSD and Frangi 
--- aim to directly capture geometric features such as line segments or tubular 
structures, which are more suited to detecting separation patterns in industrial 
surfaces. Watershed segmentation and the homography-based calibration pipeline 
serve complementary roles: the former for region separation and the latter for 
metric estimation. Compared to the deep learning methods described above, classical 
pipelines require no training data or GPU resources and offer full interpretability, 
but are generally more sensitive to parameter tuning and less robust to large 
variations in illumination or surface appearance.\\
% --------------------
\section{CREDITS}
% --------------------
\label{sec:credits}

\footnote{\url{https://www.elsevier.com/researcher/author/policies-and-guidelines/credit-author-statement}} ``CRediT (Contributor Roles Taxonomy) was introduced to recognize individual author contributions, reducing authorship disputes and facilitating collaboration''. More informally, it is a paragraph usually introduced at the end of research papers that recognizes the individual work of each author.

\begin{itemize}

\item \textbf{Pablo} contributed to the Investigation of deep learning methods 
for object detection and segmentation (YOLOv8, RT-DETRv4, Segformer++, PEM, 
Relation-DETR, ViLU-Net, Grounded-SAM, YOLO-World, Grounding DINO 1.5, and 
AnoDDPM) and participated in Writing -- Original Draft, preparing the initial 
manuscript for the corresponding sections. They also contributed to Resources, 
collecting and organizing the datasets described in the study, and to 
Visualization, preparing figures and summary tables.

\item \textbf{Jaime} contributed to the Investigation of advanced deep learning 
methods (SAM 2, EfficientSAM, YOLOv9, YOLOv8-HDSA, AnomalyDINO, CLIP-DINOv2, 
SAM2-UNet, semi-supervised defect segmentation, camera calibration survey, and 
MVTec AD 2) and participated in Writing -- Original Draft for the corresponding 
sections. Additionally, they served as the \textbf{Project POC}, acting as the 
main point of contact with the professor for doubts and task uploads. They also 
contributed to Resources and Visualization alongside the rest of the group.

\item \textbf{Miguel} contributed to the Investigation of classical computer 
vision methods (modified watershed, MAMDD, HOG+LBP+Entropy+SVM, industrial 
inspection pipelines, Harris+Bezier+Delaunay, fuzzy segmentation, SPORG-RBVS, 
ridge segmentation, subpixel edge detection, and edge detection review) and 
participated in Writing -- Original Draft for the corresponding sections. They 
also contributed to Resources and Visualization alongside the rest of the group.

\item \textbf{David} contributed to the Investigation of classical methods with 
industrial and calibration focus (HOG+SVM/PCA+RF, adaptive edge finishing, 
HOG/LBP/SIFT+KNN, grid-HOG-PCA-LightGBM, improved Canny, homography 
decomposition, camera calibration, Canny+Hu+SVM, and additional classical 
pipelines) and participated in Writing -- Original Draft for the corresponding 
sections. They also contributed to Resources and Visualization alongside the 
rest of the group.

\end{itemize}

All four members equally contributed in writing, editing, critically and 
reviewing the manuscript, and shared supervision responsibilities, 
providing oversight and leadership throughout the research activities.


% %%%%%%%%%%%%%%%%%%%%%%%%%%%%%%%%%%%%%%%%%%%%%%%%%%%%%%%%%%
% %%%%%%%%%%%%%%%%%%%%%%%%%%%%%%%%%%%%%%%%%%%%%%%%%%%%%%%%%%
% REFERENCES SECTION
% %%%%%%%%%%%%%%%%%%%%%%%%%%%%%%%%%%%%%%%%%%%%%%%%%%%%%%%%%%
% %%%%%%%%%%%%%%%%%%%%%%%%%%%%%%%%%%%%%%%%%%%%%%%%%%%%%%%%%%
\medskip

\bibliography{references.bib} 

\newpage


\end{document}